Dans ce chapitre, on se place dans $\K = \R$ ou $\C$ et $(E, N)$ est un espace vectoriel normé.

\section{Topologie d'un espace vectoriel normé}

\subsection{Ouverts }

\begin{definition}{Partie ouverte }{}
    Soit $A \subset E$. On dit que $A$ est ouverte ou que $A$ est un ouvert ssi 
    $$ \forall a \in A, \, \exists R > 0, \, B(a, R) \subset A.$$
\end{definition}

\begin{propositionnt}{}{}
    Toute boule ouverte est un ouvert.
\end{propositionnt}

\idee{
    Prendre un point de B, et trouver le rayon pour avoir une boule entièrement dans B.
}

\begin{proof}
    Soit $\omega \in E$, et $R > 0$. 

    Notons $A = B(\omega, R)$.

    Soit $a \in A$, et posons $R' = \frac{R - d(a, \omega)}{2}$. 

    Puisque $a \in A$, on a $R' > 0$.

    Soit $b \in B(a, R')$. On a : 

    \begin{align*}
        d(b, \omega) & \leqslant d(b, a) + d(a, \omega) \\
        & \leqslant R' + d(a, \omega) \\
        & \leqslant \frac{R - d(a, \omega)}{2 } + d(a, \omega) \\
        & \leqslant \frac{R + d(a, \omega)}{2} \\
        & < R \text{ car } d(a, \omega) < R.
    \end{align*}

\end{proof}

\begin{exemple}
    Dans $\R$, si $a < b$, $]a, b[ = B\left(\frac{a+b}{2}, \frac{b-a}{2}\right)$ est une boule 
    ouverte, donc un ouvert de $\R$.
\end{exemple}

\subsection{Stabilité}

\marginenote{Remarque}{Le caractère fini de l'intersection évite des contre-exemples du type $]0, \frac 1 n[$, 
dont l'intersection infinie est $\{0\}$.}

\lvspace \avspace \begin{proposition}{Opérations sur les ouverts}{}
    \begin{itemize}
        \item Toute réunion d'ouverts est un ouvert.
        \item Toute intersection \textbf{finie} d'ouverts est un ouvert.
        \item Tout produit cartésien \textbf{fini} d'ouverts est un ouvert.
    \end{itemize}
\end{proposition}

\restaurergeometrie

\idee{
    1. utiliser la def 

    2. idem


}


\begin{proof}
    \uline{Réunion :}

    \svspace Soit $I$ un ensemble non vide tel que $\forall i \in I, \, A_i$ est un ouvert de $E$, et 
    $\displaystyle B = \bigcup_{i \in I} A_i$.

    Soit $x \in B$. Il existe $i \in I$ tel que $x \in A_i$.

    Or $A_i$ est un ouvert, donc il existe $R > 0$ tel que $B(x, R) \subset A_i \subset B$, 
    donc $B$ est un ouvert.

    \avspace \qquad \uline{Intersection : }

    \svspace Soit $x \in B$. 

    Alors $\forall i \in \ninter 1 n, \, x \in A_i$. 

    $A_i$ étant ouvert, il existe $R_i > 0$ tel que $B(x, R_i) \subset A_i$.

    Posons $R = \min(R_1, \dots, R_n)$.

    Soit $y \in B(x, R)$. Alors $\forall i \in \ninter 1 n, \, y \in B(x, R_i)$, donc $y \in A_i$.

    Donc $y \in B$, et $B(x, R) \subset B$. Donc $B$ est un ouvert.

    \avspace \qquad \uline{Produit cartésien : }

    \svspace Soit $(E_1, N_1), \dots, (E_p, N_p)$ des espaces vectoriels normés, et 
    $A_1 \subset E_1, \dots, A_p \subset E_p$ des ouverts.

    Prenons $E = E_1 \times \dots \times E_p$ l'espace vectoriel normé produit muni de $N_\infty$
    et $B = A_1 \times \dots \times A_p \subset E$.

    Montrons que $B$ est un ouvert de $E$.

    Soit $x \in B$. Alors $x = (x_1, \dots, x_p)$ avec $x_i \in A_i$, où les $A_i$ sont des ouverts.

    Donc $\exists R_j > 0, \, B(x_j, R_j) \subset A_j$.

    Posons $R = \min(R_1, \dots, R_p)$.

    Soit $y \in B(x, R)$. 

    Notons $y = (y_1, \dots, y_p)$, et $N_\infty(x - y) < R$.

    Donc $\forall j \in \ninter 1 p, \, N_j(x_j - y_j) \leqslant R \leqslant R_j$, donc $y_j \in B(x_j, R_j)$.

    Donc $y_j \in A_j$, donc $y \in B$.

\end{proof}


\subsubsection{Voisinages}

\begin{definition}{Voisinage}{}
    Soit $a \in E$ et $A \subset E$. On dit que $A$ est un voisinage de $a$ ssi
    $$ \exists R > 0, \, B(a, R) \subset A$$
\end{definition}

\begin{remarque}
    En particulier, comme $a \in B(a, R)$, si $A$ est un voisinage de $a$, alors $a \in A$.
\end{remarque}


\begin{exemple}
    $\R^+$ est un voisinage de $3$ car $[2, 4] \subset \R^+$.
\end{exemple}

\begin{proposition}{Caractérisation des ouverts avec voisinages}{}
    Soit $A \subset E$. Alors $A$ est un ouvert ssi $A$ est un voisinage de chacun de ses points.
\end{proposition}

\begin{proposition}{Voisinages de l'infini}{}
    Soient $E = \R, \, A \subset \R$. On dit que $A$ est un voisinage de $+\infty$
    ssi $\exists c \in \R, \, [c, +\infty[ \subset A$. 

    On dit que $A$ est un voisinage de $-\infty$ ssi 
    $\exists c \in \R, \, ]-\infty, c] \subset A$.
\end{proposition}


\subsection{Fermés}

\begin{definition}{Fermé}{}
    Soit $A \subset E$. On dit que $A$ est un fermé ssi son complémentaire 
    $E \setminus A = C_E B$ est un ouvert.
\end{definition}


\begin{exemple}
    $\emptyset$ et $E$ sont des fermés de $E$. 
\end{exemple}

\begin{theoreme}{Caractérisation séquentielle des fermés}{}

    Soit $B \subset E$. 

    $B$ est un fermé ssi pour toute suite $(u_n)_{n \in \N}$ d'éléments de $B$ qui converge 
    vers $\ell \in E$, on a $\ell \in B$.

\end{theoreme}

\idee{
    $\Rightarrow$ : par l'absurde, on suppose la limite dans le complémentaire. 
}

\begin{proof}
    Montrons $\Rightarrow$. Supposons donc que $B$ est un fermé.

    Soit $(u_n)$ une suite d'éléments de $B$ qui converge vers $\ell \in E$.

    Montrons que $\ell \in B$.

    Supposons $\ell \in A$. $B$ est fermé, donc $A$ est ouvert, donc $\exists R > 0, \, B(\ell, R) \subset A$.

    Or $u_n \limit n \infty \ell$. Prenons $\varepsilon = \frac R 2$. 

    Alors $\exists n_0 \in \N, \, \forall n \geqslant n_0, \, d(u_{n}, \ell) \leqslant \varepsilon$. 

    En particulier, $d(u_{n_0}, \ell) \leqslant \frac R 2 < R$, donc $u_{n_0} \in A$, or $u_{n_0} \in B$, ce qui est absurde.

    Donc $\ell \not \in A$, donc $\ell \in B$.

    \avspace $\Leftarrow$ Supposons maintenant que $B$ n'est pas fermé. Donc $A$ n'est pas ouvert, ce qui s'écrit : 

    $$\exists a \in A, \, \forall R > 0, \, \exists x \in B(a, R) \text{ tel que } x \not \in A$$

    En particulier, pour $n \in \N^*$, $R = \frac 1 n$ fournit un $x_n \in B(a, \frac 1 n)$ tel que $x_n \not \in A$, donc $x_n \in B$ 
    et $d(x_n, a) < \frac 1 n$.

    Donc $(x_n)$ est une suite de $B$ qui converge vers $a \in A$, avec $a \not \in B$.

    C'est la contraposée de la réciproque, d'où l'équivalence. 
\end{proof}


\begin{proposition}{Fermés usuels}{}
    Une boule fermée, une sphère, un singleton sont des fermés. 
\end{proposition}

\idee{
    Caractérisaiton séquentielle avec distance au centre.
}

\begin{proof}
    Soit $R \geqslant 0$, et $\omega \in E$.

    Considérons la boule fermée $\overline{B}(\omega, R)$.

    \begin{itemize}
        \item Si $R > 0$, alors c'est une boule fermée
        \item Si $R = 0$, c'est un singleton. 
    \end{itemize}

    Soit $(x_n)_{n \in \N}$ suite de $B$ qui converge vers $\ell \in E$.

    $\forall n \in \N, \, N(x_n - \omega) \leqslant R$.

    Or $x_n \limit n \infty \ell$, donc $x_n - \omega \limit n \infty \ell - \omega$.

    Donc $N(x_n - \omega) \limit n \infty N(\ell - \omega)$.

    Donc par passage à la limite dans $\R$, $N(\ell - \omega ) \leqslant R$, donc $\ell \in B$.

    Donc $B$ est un fermé par caractérisation séquentielle.

    Pour la sphère, la preuve est identique, en remplaçant $\leqslant$ par $=$.
\end{proof}


\subsubsection{Stabilité }

\begin{proposition}{Opérations sur les fermés}{}
    \begin{itemize}
        \item Toute réunion \textbf{finie} de fermés est un fermé.
        \item Toute intersection quelconque de fermés est un fermé.
        \item Tout produit cartésien \textbf{fini} de fermés est un fermé.
    \end{itemize}
\end{proposition}

Astuce pour ne pas inverser avec les ouverts : se rappeler que chaque ensemble est l'union des singletons 
de ses éléments, or tout ensemble n'est pas fermé. 

\idee{
    Utiliser la définition de fermé avec les ouverts. 

    Pour le produit cartésien, utiliser la caractérisation séquentielle.
}

\begin{proof}
    Soit $B_1, \dots, B_n$ des fermés de $E$, et $\displaystyle C = \bigcup_{i=1}^n B_i$.

    Ainsi, $E \setminus C = \underbrace{\bigcap_{i=1}^n \underbrace{(E \setminus B_i)}_{\text{ouverts}}}_{\text{ouverts}}$.

    Donc $C$ est un fermé.

    \lvspace Soit $I$ un ensemble non vide, et $(B_i)_{i \in I}$ une famille de fermés de $E$.

    Posons $\displaystyle C = \bigcap_{i \in I} B_i$.

    Ainsi, $\displaystyle E \setminus C = \bigcup_{i \in I} (E \setminus B_i)$, qui est une réunion d'ouverts, donc un ouvert.

    Donc $C$ est un fermé.

    \lvspace Soit $(E_1, N_1), \dots, (E_p, N_p)$ des espaces vectoriels normés, et
    $E = E_1 \times \dots \times E_p$ l'espace vectoriel normé produit muni de $N_\infty$.

    $\forall j \in \ninter 1 p, \, B_j$ est un fermé de $E_j$.

    Soit $B = B_1 \times \dots \times B_p$. 

    Soit $(x_n)_{n \in \N}$ une suite de $B$. Notons $\forall n \in \N, \, x_n = (x_{n, 1}, \dots, x_{n, p})$.

    Supposons que $(x_n)$ converge vers $\ell \in E$. 

    Notons $\ell = (\ell_1, \dots, \ell_p)$.

    On a $\forall j \in \ninter 1 p, \, x_{n, j} \limit n \infty \ell_j$.

    Or $\forall n \in \N, \, x_{n, j} \in B_j$, et $B_j$ est un fermé, donc $\ell_j \in B_j$.

    Donc $\ell \in B$, et $B$ est un fermé par caractérisation séquentielle.

\end{proof}


\qquad \uline{Conséquence : }

\svspace Tout ensemble fini est fermé, car réunion finie de singletons.

\begin{theoreme}{Hors programme}{}
    Les seules parties à la fois ouvertes et fermées de $E$ sont $\emptyset$ et $E$.
\end{theoreme}

\idee{
    Par l'absurde, on prend deux ensembles ouverts et fermés complémentaires, et 
    on cherche à aller sur la frontière. 
}

\begin{proof}
    Supposons par l'absurde qu'il existe $A \subset E$ telle que $A \neq \emptyset$, $A \neq E$, et $A$ est à la fois ouvert et fermé.

    Alors $B = E \setminus A$ n'est pas vide, et est ouvert et fermé. 

    Puisque $A \neq \emptyset, \, B \neq \emptyset$, il existe $a \in A$, et $b \in B$, 
    et considérons 
    
    $S = \left\{\lambda b + (1 - \lambda )a, \, \lambda \in \inter 0 1\right\}$

    Notons $C = \left\{\lambda \in \inter 0 1, \, \lambda b + (1 - \lambda )a \in A\right\}$

    Donc $0 \in C$, donc $C \neq \emptyset$.

    De plus, $C \subset \R$ non vide et borné. Posons alors $\lambda_0 = \sup C \in \R$, et 
    $d = \lambda_0 b + (1 - \lambda_0)a \in E$.

    On a $\forall t \in C, \, t \leqslant \lambda_0$

    Posons pour $\displaystyle n \in \N$, $\displaystyle t_n = \lambda_0 + \frac{1 - \lambda_0}{n + 1}$. 

    On a $\forall t_n \in \interoc{\lambda_0}{1}$, donc $t_n \not \in C$, donc $t_n b + (1 - t_n)a \in B$.

    Or $t_n b + (1 - t_n)a \limit n \infty d$. Or $B$ est fermé, donc $d \in B$.

    Par caractérisation séquantielle du sup, il existe une suite $(u_n)_{n \in \N}$ de $C$ qui converge 
    vers $\lambda_0$.

    Donc $u_n b + (1 - u_n)a \in A$, or $u_n b + (1 - u_n)a \limit n \infty d$.

    Or $A$ est fermé, donc $d \in A$.

    Donc $A \cap B \neq \emptyset$ absurde.
\end{proof}

\begin{theoreme}{Nature sous espace vectoriel de dimension finie}{}
    Tout sous espace vectoriel de dimension finie de $E$ est un fermé. 
\end{theoreme}

\idee{
    Utilisation de Bolzano-Weierstrass et unicité de la limite. 
}

\begin{proof}
    Soit $F$ un sous espace vectoriel de dimension finie de $E$. 

    \begin{itemize}
        \item Si $F = \{0_E\}$, c'est un singleton, donc un fermé.
        \item Sinon, supposons $F \neq \emptyset$. 
    \end{itemize}

    Soit $(u_n)_{n \in \N}$ une suite de $F$ qui converge vers $\ell \in E$.

    Cette suite converge, donc c'est une suite bornée de $F$ pour $N$. 

    Donc par le théorème de Bolzano-Weierstrass, on peut en extraite une sous-suite $(u_{\varphi(n)})_{n \in \N}$ qui converge vers
    $\ell_1 \in F$.

    Or $u_n \limit n \infty \ell$, donc $\ell_1 = \ell$, donc $\ell \in F$.
\end{proof}

Montrons qu'au contraire, cela n'est pas vrai en dimension infinie.

\begin{proof}
    $E = \R [X]$ muni de $N: P \mapsto \int_{0}^1 |P(t)| dt$.

    Soit $H = \{P \in E, \, P(1) = 0\}$, qui est un sous espace vectoriel de $E$.

    Montrons que $H$ n'est pas fermé.

    Posons $\forall n \in \N, \, P_n = X^n - 1$. On a donc bien $P_n \in H$.

    Posons $Q = -1$. Alors 

    $$ N(p_n - Q ) = \int_0^1 |X^n| dt = \frac 1 {n+1} \limit n \infty 0.$$

    Donc $P_n \limit n \infty Q$ et pourtant $Q \not \in H$.

    Donc $H$ n'est pas fermé.
\end{proof}


\subsection{Intérieur, adhérence, frontière}

\begin{definition}{Point intérieur}{}
    Soit $A \subset E$, et $a \in E$.

    On dit que $a$ est un point intérieur à $A$ ssi $\exists R > 0, \, B(a, R) \subset A$,
    ssi $A$ est un voisinage de $a$. 
\end{definition}



\begin{definition}{Intérieur d'un ensemble}{}
    Soit $A \subset E$. 

    On appelle intérieur de $A$, et on note $\mathring{A}$, l'ensemble des points intérieurs à $A$, 
    c'est à dire : 

    $$ \mathring{A} = \{a \in A, \, \exists R > 0, \, B(a, R) \subset A\}.$$
\end{definition}

\begin{proposition}{Inclusion de l'intérieur}{}
    On a $\mathring{A} \subset A$ avec égalité ssi $A$ est un ouvert.
\end{proposition}

\begin{proof}
    Par définition, $\mathring{A} \subset A$.

    Si $A$ est un ouvert, alors $A$ est un voisinage de chacun de ses points, donc $\mathring{A} = A$.

    Réciproquement, si $\mathring{A} = A$, 
    alors $a \in \mathring A$ donc par def, $\exists R > 0, \, B(a, R) \subset A$, donc $A$ est un ouvert.
\end{proof}

\begin{corollaire}{intérieur d'une boule fermée}{}
    Soit $\omega \in E, \, R_0 > 0$, et $A = B'(\omega, R_0)$, alors $\mathring{A} = B(\omega, R_0)$.
\end{corollaire}

Je ne garantis pas que l'une ou l'autre des preuves soit juste (même si celle du cours 
doit l'être à peu près). 

\idee{ Montrer l'inclusion, puis que les éléments de la sphère ne sont pas 
intérieurs avec une suite. }

\begin{proof}

    \uline{Preuve donnée en cours :}
    Soit $a \in B(\omega, R_0)$. 

    On sait que $B(\omega, R_0)$ est un ouvert. 

    Donc $\exists R > 0, \, B(a, R) \subset B(\omega, R_0)$ 
    donc $B(a, R) \subset A$ donc $a \in \mathring{A}$.

    Ainsi, $B(\omega, R_0) \subset \mathring{A} \subset B'(\omega, R_0)$.

    Soit $x \in B'(\omega, R_0) \setminus B(\omega, R_0)$.

    Alors $x \in S(\omega, R_0)$. 

    Montrons que $x \not \in \mathring A$. 

    Soit $R > 0$. Posons $u = \frac{x - \omega}{N(x - \omega)}$ et $y = x + \frac R 2 u$.

    Alors $d(y, x) = N(y - x) = N\left( \frac R 2 u \right) = \frac R 2 < R$, donc $y \in B(x, R)$.
    et $d(\omega, y) = N(y - \omega)$, 

    \begin{align*}
        N(y - \omega) & = N\left(x + \frac R 2 \frac{x - \omega}{N(x - \omega)} - \omega\right) \\
        & = N\left( \left( N(x - \omega) + \frac R 2 \right) \left( \frac{x - \omega}{N(x - \omega)}\right) \right) \\
        & = N(x - \omega) + \frac R 2 \\
        & > R_0.
    \end{align*}

    Donc $y \not \in B(\omega, R_0)$, donc $x \not \in \mathring A$.

    Donc $\mathring A = B(\omega, R_0)$

    \uindent{Preuve alternative }

    Si \( x \in B(\omega, R_0) \), alors \( d(x, \omega) < R_0 \).
    Posons \( \varepsilon = R_0 - d(x, \omega) > 0 \).
    La boule ouverte \( B(x, \varepsilon) \) est contenue dans \( B'(\omega, R_0) \).
    Donc \( x \) est un point intérieur de \( A \).
    Donc \( B(\omega, R_0) \subset \mathring{A} \).

    Si \( x \in \mathring{A} \), alors il existe \( \varepsilon > 0 \) tel que
    \( B(x, \varepsilon) \subset B'(\omega, R_0) \).
    Donc pour tout \( y \in B(x, \varepsilon) \), \( d(y, \omega) \leq R_0 \).
    En particulier, \( d(x, \omega) < R_0 \) (strict, sinon il n'y aurait pas de marge pour contenir une boule autour de \( x \)).
    Donc \( x \in B(\omega, R_0) \).
    Donc \( \mathring{A} \subset B(\omega, R_0) \).

    Conclusion :
    \( \mathring{A} = B(\omega, R_0) \).
\end{proof}


\begin{theoreme}{Intérieur comme réunion d'ouverts}{}
    Soit $A \subset E$.

    \begin{itemize}
        \item Alors $\mathring{A}$ est la réunion des ouverts contenus dans $A$.
        \item C'est le plus grand ouvert dans $A$.
    \end{itemize}
\end{theoreme}

\begin{proof}
    Notons $(U_i)_{i \in I}$ la famille des ouverts contenus dans $A$, 
    
    et $C = \bigcup_{i \in I} U_i$.

    $\mathring A$ est un ouvert inclus dans $A$, donc c'est l'un des $U_i$, donc $\mathring A \subset C$.

    Soit $x \in C$. Alors $\exists i \in I, \, x \in U_i$.

    Or $U_i$ est un ouvert, donc $\exists R > 0, \, B(x, R) \subset U_i$. 

    Mais alors $B(x, R) \subset A$. Donc $x \in \mathring A$.

    Donc $C \subset \mathring A$.

    Ainsi $\mathring A$ est un ouvert inclus dans $A$. Si $B$ est un ouvert inclus dans 
    $A$, c'est l'un des $U_i$ donc $B \subset C = \mathring A$.

    Donc $\mathring A$ est le plus grand ouvert inclus dans $A$.
\end{proof}

\begin{definition}{Point adhérent}{}
    Soit $A \subset E$ et $x \in E$.

    On dit que $x$ est adhérent à $A$ ssi 

    $$ \forall \varepsilon > 0, \, B(x, \varepsilon) \cap A \neq \emptyset.$$
\end{definition}

\begin{theoreme}{Caractérisation séquentielle des points adhérents}{}
    Soit $A \subset E$ et $x \in E$.

    Alors $x$ est adhérent à $A$ ssi il existe une suite $(a_n)_{n \in \N}$ d'éléments de $A$ 
    qui converge vers $x$.
\end{theoreme}

\begin{proof}
    Soit $x \in E$ adhérent à $A$. 

    Alors $\forall R > 0, \, \exists a \in A, \, d(x, a) < R$.

    Soit $n \in \N*$, avec $R = \frac 1 n$n cela fournit $a_n \in A$ tel que $d(x, a_n) < \frac 1 n$.

    Donc $a_n \limit n \infty x$.

    Réciproquement, soit $x \in E$ tel qu'il existe une suite $(a_n)_{n \in \N}$ d'éléments 
    de $A$ qui converge vers $x$. 

    Soit $R > 0$ 

    Par def de la limite, $\exists n_0 \in \N, \, \forall n \geqslant n_0, \, d(x, a_n) < R$.

    Donc $a_{n_0} \in B(x, R) \cap A$. 
\end{proof}

\begin{definition}{Adhérence d'un ensemble}{}
    Soit $A \subset E$. 
    L'ensemble des points adhérents à $A$ est appelé adhérence de $A$, et noté 
    $\overline{A}$.
\end{definition}

\begin{proposition}{Inclusion de l'adhérence}{}
    Soit $A \subset E$. Alors $A \subset \overline{A}$ avec égalité ssi $A$ est un fermé.
\end{proposition}

\begin{proof}
    Soit $x \in A$. Pour tout $R > 0, \, x \in B(x, R) \cap A$ donc $x \in \overline{A}$.
    Donc $A \subset \overline{A}$.

    \avspace Si $A$ est fermé. 

    Soit $y \in \overline{A}$. Alors il existe une suite $(a_n)_{n \in \N}$ d'éléments
    de $A$ qui converge vers $y$.

    Or $A$ est fermé, donc $y \in A$, et $\overline{A} \subset A$ et donc $\overline{A} = A$.

    \qquad \avspace \uline{Si $\overline{A} = A$ : }

    \svspace Soit $(a_n)_{n \in \N}$ une suite d'éléments de $A$ qui converge vers $\ell \in E$.

    Alors par caractérisation séquentielle, $\ell \in \overline{A} $ Donc $A$ fermé. 
\end{proof}

\begin{proposition}{Adhérence des boules}{}
    Soit $\omega \in E$ et $R > 0$. Alors
     $$\overline{B(\omega, R)} = B'(\omega, R)$$
\end{proposition}

\begin{proof}
    Soit $x \in \overline{B(\omega, R)}$.

    Par caractérisation séquentielle, il existe une suite $(a_n)_{n \in \N}$ d'éléments de
    $B(\omega, R)$ qui converge vers $x$.

    $N(a_n - \omega) < R$, donc par passage à la limite, $N(x - \omega) \leqslant R$. 

    Donc $x \in B'(\omega, R)$, et $\overline{B(\omega, R)} \subset B'(\omega, R)$.

    \avspace \qquad \uline{Réciproquement : }

    \svspace Soit $x \in B'(\omega, R)$.

    Si $x = \omega$, alors $x \in B(\omega, R)$ donc $x \in \overline{B(\omega, R)}$.

    Si $x \neq \omega$, posons $u = \frac{x - \omega }{N(x - \omega)}$.

    Posons $a_n = \omega + \frac n {n + 1} N(x - \omega)u$. 

    \begin{align*}
        N(a_n - \omega) & = N\left( \frac n {n + 1} N(x - \omega) u \right) \\
        & = \frac n {n + 1} N(x - \omega) 
        & \leqslant \frac n {n + 1} R < R.  
    \end{align*}

    Donc $a_n \in B(\omega, R)$.

    \begin{align*}
        N(a_n - x) & = N\left( \omega - x + \frac n {n + 1} (x - \omega) \right) \\
        & = N\left( \frac 1 {n + 1}(\omega - x) \right) \\
        & = \frac 1 {n + 1} N(\omega - x) \limit n \infty 0.
    \end{align*}

    Donc $a_n \limit n \infty x$, 

    Donc $x \in \overline{B(\omega, R)}$, donc $B'(\omega, R) \subset \overline{B(\omega, R)}$.
\end{proof}

\begin{theoreme}{Adhérence comme intersection de fermés}{}
    Soit $A \subset E$.
    $\overline{A}$ est l'intersection de tous les fermés contenant $A$.
    C'est le plus petit fermé contenant $A$.
\end{theoreme}

\begin{proof}
    Notons $(F_i)_{i \in I}$ la famille des fermés contenant $A$, et
    $ \displaystyle C = \bigcap_{i \in I} F_i$. Montrons que $C = \overline{A}$.

    $\overline{A}$ est un fermé contenant $A$, donc c'est l'un des $F_i$, donc $ C \subset \overline{A}$.

    Soit $x \in \overline{A}$, et soit $i \in I$. 

    Par caractérisation séquentielle de l'adhérence, il existe une suite $(a_n)_{n \in \N}$ d'éléments de $A$ qui converge vers $x$.

    Mais $A \subset F_i$, donc $a_n \in F_i$.

    Or $F_i$ fermé, donc $x \in F_i$ et donc $\overline A \subset F_i$, 
    et alors $\overline A \subset C$.

    $\overline A$ est un fermé contenant $A$. 

    Soit $B$ un fermé contenant $A$. C'est l'un des $F_i$, donc $C \subset B$, 
    donc $\overline A \subset B$.
\end{proof}


\uindent{Bilan}

\avspace Si $A \subset E$, alors
$$\mathring A \subset A \subset \overline A$$


\begin{proposition}{Hors programme}{}
    Soit $A \subset E, \, B \subset E$. Alors 
    \begin{itemize}
        \item $\overline{E \setminus A} = E \setminus \mathring{A}$
        \item $\overline{A \cup B} = \overline{A} \cup \overline{B}$
        \item $\overline{A \cap B} \subset \overline{A} \cap \overline{B}$
    \end{itemize} 
\end{proposition}


\begin{proof}
    On a $\mathring A \subset A$. 

    Donc $E \setminus A \subset \underbrace{E \setminus \underbrace{\mathring A}_{\text{ouvert}}}_{\text{fermé}}$.

    Donc $\overline{E \setminus A} \subset E \setminus \mathring A$.


    \uindent{Réciproquement }

    \svspace $E \setminus A \subset \overline{E \setminus A}$

    Donc $E \setminus \overline{E \setminus A} \subset A$.

    Donc $E \setminus \overline{E \setminus A} \subset \mathring A$ 

    Donc $E \setminus \mathring A \subset \overline{E \setminus A}$. 

    \uindent {Deuxième point }

    \svspace $A \subset \overline A$, et $B \subset \overline B$,

    Donc $A \cup B \subset \underbrace{\overline A \cup \overline B}_{\text{fermé}}$.

    Donc $\overline{A \cup B} \subset \overline A \cup \overline B$.


    \uindent{Réciproquement }

    \svspace $A \subset A \cup B$, donc $A \subset \overline{A \cup B}$.

    Donc $\overline A \subset \overline{A \cup B}$.

    De même, $\overline B \subset \overline{A \cup B}$.

    Donc $\overline A \cup \overline B \subset \overline{A \cup B}$.

    \uindent{Troisième point }

    $A \cap B \subset A$, donc $A \cap B \subset \overline A$. 

    Donc $\overline {A \cap B} \subset \overline A$.

    De même $\overline {A \cap B} \subset \overline B$.

    Donc $\overline {A \cap B} \subset \overline A \cap \overline B$.
\end{proof}

\textbf{Attention}, l'inclusion réciproque n'est pas vraie : 

Avec $E = \R$, $A = \interoo 0 1$, $B = \interoo 1 2$, on a 
$\overline A = [0, 1]$, $\overline B = [1, 2]$, donc $A \cap B = \emptyset$, et 
$\overline A \cap \overline B = \{1\}$.

\begin{definition}{Frontière d'un ensemble}{}
    Soit $A \subset E$. 

    On appelle frontière de $A$ et on note $\operatorname{Fr}(A)$ la partie de $E$ définie 
    par 
    $$ \operatorname{Fr}(A) = \overline A \setminus \mathring A$$
\end{definition}

\begin{exemple}
    Si $A = \Q$, 

    $\mathring A = \emptyset$, $\overline A = \R$, donc $\operatorname{Fr}(A) = \R$.
\end{exemple}


\begin{proposition}{Nature de la frontière}{}
    Soit $A \subset E$. Alors $\Fr(A)$ est un fermé et $\Fr(A) = \Fr(E \setminus A)$.
\end{proposition}

\begin{proof}
    $\Fr A = \overline A \cap (E \setminus \mathring A)$. 

    De plus $E \setminus \mathring A = \overline{E \setminus A}$.

    Donc $\Fr A = \overline A \cap \overline{E \setminus A} = \Fr (E \setminus A)$ 
\end{proof}

\subsection{Partie dense }

\begin{definitionnt}{Partie dense inclusion}{}
    Soient $A, B \in \mathscr P(E)$.

    On dit que $A$ est dense dans $B$ ssi 
    $$A \subset B \subset \overline A$$
\end{definitionnt}

\begin{definitionnt}{Partie dense égalité}{}
    Soit $A \in \mathscr P(E)$.

    On dit que $A$ est dense dans $E$ ssi $\overline A = E$.
\end{definitionnt}

\begin{exemple}
    \begin{enumerate}
        \item $\Q$ est dense dans $\R$.
        \item $\operatorname{GL}_p(\K)$ est dense dans $\mathcal M_p(\K)$.
        \item Si $H$ est un hyperplan de $E$, alors $H$ est fermé ou dense
    \end{enumerate}


    \uindent{Preuve du 2. }

    \svspace Soit $A \in \mathcal M_p(\K)$.

    Pour $n \in \N^*$, posons $B_n = A - \frac 1 n I_p$.

    Par théorème des valeurs propres, $B_n \limit n {+ \infty} A$.

    Or 
    \begin{align*}
        \det(B_n) & = \det((A_{i, j} - \frac 1 n \delta_{i, j})_{1 \leqslant i, j \leqslant p}) \\
        & = \sum_{ \sigma \in S_p} \varepsilon(\sigma) \prod_{i=1}^p (A_{i, \sigma(i)} - \frac 1 n \delta_{i, \sigma(i)}) 
    \end{align*}

    Posons $x = \frac 1 n > 0$. 

    Donc $\det(B_n) $ est un polynôme en $x$ de degré inférieur à $p$. 

    Le coef de $x^p$ provient unisuement de $\sigma = \text{Id}$, et vaut $(-1)^p$.

    Ce polynôme a au plus $p$ racines. 

    D'où $\exists n_0 \in \N, \, \forall n \geqslant n_0, \, \frac 1 n$ n'en est pas une racine. 

    Donc $\forall n \geqslant n_0, \, B_n \in \operatorname{GL}_p(\K)$.


    \uindent{Preuve du 3. }

    \svspace Soit $H$ un hyperplan de $E$. Si $\varphi$ est une forme linéaire non nulle
    telle que $H = \ker \varphi$.

    Soit $a \in E $tel que $\varphi(a) \neq 0$, alors $E = H \oplus \operatorname{Vect}(a)$.

    En effet, pour $x \in E$ : 
    
    \uindent{Analyse }

    \svspace Supposons que $x = h + \lambda a$ avec $h \in H$ et $\lambda \in \K$.

    Alors $\varphi(x) = \varphi(h) + \lambda \varphi(a) = \lambda \varphi(a)$, donc
    $\lambda = \frac{\varphi(x)}{\varphi(a)}$ et $h = x - \frac{\varphi(x)}{\varphi(a)} a$.

    d'où l'unicité sous réserve d'existence. 

    \uindent{Synthèse }

    Posons $\lambda = \frac{\varphi(x)}{\varphi(a)}$ et $h = x - \lambda a$.

    Alors $x = h + \lambda a$ par calcul. $\lambda \in \K$. 

    $\varphi(h) = \varphi(x) - \frac{\varphi(x)}{\varphi(a)} \varphi(a) = 0$, donc $h \in H$.

    Supposons que $H$ n'est pas fermé, et montrons qu'il est dense.

    Il existe une suite $(h_n)_{n \in \N}$ d'éléments de $H$ telle que 
    $h_n \limit n \infty \ell$ avec $\ell \not \in H$, donc $\varphi(\ell) \neq 0$.

    Donc $E = H \oplus \operatorname{Vect}(\ell)$.

    Soit $x \in E$. Décomposons $x = h_x + \lambda_x \ell$ où $h_x \in H$ et $\lambda_x \in \K$.

    Posons $x_n = h_x + \lambda_x h_n$. On a bien $x_n \in H$, et $x_n \limit n \infty x$ 
    d'où $\overline H = E$.


\end{exemple}


\subsection{Invariance}

\begin{theoreme}{Invariance topologique en dimension finie}{}
    Les notions topologiques (ouverts, fermés, intérieur, adhérence, frontière) sont invariantes
    par remplacement de la norme par une norme équivalente. 
\end{theoreme}

\begin{proof}
    La notion de fermé qui vient de la caractérisation séquentielle est invariante car la notion de 
    limite de suire l'est. 

    La notions d'ouvert ausi par complémentarité. 

    Les autres notions également car elles s'expriment en fonction des deux premières.
\end{proof}


\begin{corollaire}{invariance topologique en dimension finie}{}
    En dimension finie, ces notions topologiques ne dépendent pas du choix de la norme. 
\end{corollaire}



\subsubsection{Topologie relative}

Dans toute cette partie, on fixe $A$ partie non vide de $E$.

\begin{definition}{Ouvert relatif}{}
    Soit $U \subset A$. On dit que $U$ est un ouvert relatif de $A$ 
    ssi $\forall a \in U, \, \exists R > 0, \, \forall x \in A, \, D(a, x) < R \Rightarrow x \in U$.

    C'est à dire $B(a, R) \cap A \subset U$.
 
\end{definition}

\begin{definition}{Voisinage relatif}{}
    Soit $U \subset A$, et $a \in A$. 

    On dit que $U$ est un voisinage relatif de $a$ vis à vis de $A$ ssi 
    $\exists R > 0, \, B(a, R) \cap A \subset U$.

    ainsi un ouvert relatif est voisinage relatif de chacun de ses points.
\end{definition}

\begin{exemple}
    $E = \R$, $A = \R^+$, $U = \interco 0 1$ est un voisinage relatif de $0$ 
    car $B(0, 1) \cap \R^+ \subset U$. 
\end{exemple}


\begin{definition}{Fermé relatif}{}
    Soit $V \subset A$. 

    On dit que $V$ est un fermé relatif de $A$ ssi $A \setminus V$ est un ouvert relatif de $A$.
\end{definition}


\begin{proposition}{Caractérisation des fermés relatifs}{}
    Soit $V \subset A$ est un fermé relatif de $A$ ssi il existe $W$ fermé de $E$ tel que $V = W \cap A$.
\end{proposition}

\idee{
    $\Rightarrow$ Utilisation de $\overline V$ et absurde. 

    $\Leftarrow$ Utilisation du complémentaire. 
}

\begin{proof}
    $\Rightarrow$ : Supposons que $V$ est un fermé relatif de $A$.

    Posons $W = \overline{V}$, $W$ est un fermé de $E$ et montrons que $V = W \cap A$.

    $V \subset A$, $V \subset \overline V = W$, donc $V \subset W \cap A$.

    Supposons que $x \in \overline V \cap A$ mais que $x \not \in V$.

    Notons $U = A \setminus V$, on a donc $x \in U$ et par def, $U$ est un ouvert relatif. 

    Donc $\exists R > 0, \, B(x, R) \cap A \subset U$.

    Mais $x \in  \overline V$ donc $B(x, R) \cap V \neq \emptyset$.

    Donc il existe un $y \in A$ tel que $y \in V$, $y \in B(x, R)$. Donc $y \in U$.

    Absurde car $U \cap V = \emptyset$.

    Donc $V = W \cap A$.

    \uindent{Réciproquement }

    Supposons qu'il existe $W$ fermé tel que $V = W \cap A$. 

    Alors $U = A \setminus V = A \setminus (A \cap W)$

    Notons $C = E \setminus W$ ouvert.

    Soit $a \in U$. Alors $a \in C$ et $C$ est ouvert. 

    Donc $\exists R > 0, B(a, R) \subset C$.

    Donc $B(a, R) \subset C \cap A = U$ 

    Donc $U$ est un ouvert relatif de $A$, donc $V$ est un fermé relatif de $A$.
\end{proof}


\begin{proposition}{Caractérisation séquentielle des fermés relatifs}{}
    Soit $V \subset A$. 

    Alors $V$ est un fermé relatif de $A$ ssi pour toute suite $(a_n)_{n \in \N}$ d'éléments de $V$
    qui converge vers $\ell \in A$, on a $\ell \in V$.
    
\end{proposition}

\idee{
    Utilisation de la caractérisation séquentielle des fermés, et proposition précédente.
}

\begin{proof}
    Supposons que $V$ est un fermé relatif de $A$. 
    Alors $V = \overline V \cap A$. 

    Soit $(a_n)_{n \in \N}$ une suite d'éléments de $V$ qui converge vers $\ell \in A$.

    Alors $\ell \in \overline V$, donc $\ell \in \overline V \cap A$ donc $\ell \in V$.

    \uindent{Réciproquement }

    \svspace Supposons que pour toute suite $(a_n)_{n \in \N}$ d'éléments de $V$ qui converge vers
    $\ell \in A$, on a $\ell \in V$, et montrons que $V = \overline V \cap A$. 

    On a $V \subset \overline V$ et $V \subset A$, donc $V \subset \overline V \cap A$.

    Soit $\ell \in \overline V \cap A$.

    Alors $\ell \in \overline V$ donc il existe une suite $(a_n)_{n \in \N}$ d'éléments de $V$ qui converge vers $\ell$.

    Or $\ell \in A$, donc par hypothèse, $\ell \in V$.

    Donc $\overline V \cap A \subset V$, donc $V = \overline V \cap A$.

    Donc $V$ est un fermé relatif de $A$.
\end{proof}


\begin{proposition}{Caractérisation des ouverts relatifs}{}
    Soit $U \subset A$. 

    Alors $U$ est un ouvert relatif de $A$ ssi il existe un ouvert $V$ de $E$ tel que 
    $U = V \cap A$.
    
\end{proposition}

\begin{proof}
    Supposons que $U$ est un ouvert relatif de $A$.

    Alors $V = a \setminus U$ est un fermé relatif de $A$.

    Donc $V = \overline V \cap A$, donc $U = A \setminus (\overline V \cap A)$. 

    Donc $U = A \cap \underbrace{(E \setminus \overline V)}_{\text{ouvert de } E}$

    \avspace Supposons qu'il existe $C$ ouvert de $E$ tel que $U = C \cap A$.

    Soit $x \in U$. Alors $x \in C$ et $C$ est ouvert.

    Donc $\exists R > 0, \, B(x, R) \subset C$.

    Donc $B(x, R) \cap A \subset U$ donc $U$ est un ouvert relatif de $A$.
\end{proof}



\section{Etude locale d'une application}

\begin{definition}{Limite}{}
    Soient $(E, N)$ et $(E', N')$ deux espaces vectoriels normés. 

    Soit $A \subset E$ non vide, et $\fonction{f}{A}{E'}{x}{f(x)}$

    Soit $a \in \overline A$, et $\ell \in E'$.

    On dit que $f$ a pour limite $\ell$ en $a$ ssi 
    $$\forall \varepsilon > 0, \, \exists \eta > 0, \, \forall x \in A, \, N(x - a) < \eta \Rightarrow N'(f(x) - \ell) < \varepsilon$$

    c'est à dire $d(x, a) \leqslant \eta \Rightarrow d(f(x), \ell) \leqslant \varepsilon$.

    c'est à dire $x \in B'(a, \eta) \Rightarrow f(x) \in B'(\ell, \varepsilon)$.

    $\ell$ est alors unique, on l'appelle limite de $f$ en $a$ et on note 

    $$f(x) \xrightarrow[x \to a ]{N'} \ell$$

\end{definition}


\begin{proof}
    Preuve de l'unicité 

    Supposons que $\ell_1$ et $\ell_2$ conviennent. Soit $\varepsilon > 0$.

    Par def, $\exists M_1 > 0, \, \forall x \in A, \, d(x, a) < M_1  \Rightarrow d(f(x), \ell_1) \leqslant \varepsilon$

    $\exists M_2 > 0, \, \forall x \in A, \, d(x, a) < M_2  \Rightarrow d(f(x), \ell_2) \leqslant \varepsilon$

    Posons $M_3 = \min(M_1, M_2)$, et soit $x \in A$ tel que $d(x, a) < M_3$.

    Un tel $x$ existe car $a \in \overline A$.

    On a $d(\ell_1, \ell_2) \leqslant d(\ell_1, f(x)) + d(f(x), \ell_2) \leqslant 2 \varepsilon$.

    Vrai pour tout $\varepsilon > 0$, donc $d(\ell_1, \ell_2) = 0$, donc $\ell_1 = \ell_2$.
\end{proof}

\begin{remarque}
    Cette définiton est invariante si on remplace $N$ et $N'$ par des normes équivalentes.
\end{remarque}

\begin{proof}
    Soit $||\cdot||$ une norme équivalente à $N$, norme de $E$, 
    et $||\cdot||'$ une norme équivalente à $N'$, norme de $E'$.

    Alors $\exists \alpha, \beta, \gamma, \tau > 0, \forall x \in E, \, \beta ||x|| \leqslant N(x) \leqslant \alpha ||x||$.
    et $\forall y \in E', \, \gamma ||y||' \leqslant N'(y) \leqslant \tau ||y||'$.

    Supposons que $f(x) \limitnorme{N'} x a \ell$. Ainsi, $\forall \varepsilon > 0, \exists \eta > 0, \, \forall x \in A, \, N(x - a) < \eta \Rightarrow N'(f(x) - \ell) < \varepsilon$.

    Soit $\varepsilon > 0$, et posons $\eta' = \frac \eta \alpha$ où $\eta$ vient de la 
    définitions de la limite avec $\gamma \varepsilon$ das le rôle de $\varepsilon$.

    Soit $x \in A$. Supposons que $||x - a|| < \eta'$. 

    Alors $||f(x) - \ell||' < \frac 1 \gamma N(f(x) - \ell)$

    Or \begin{align*}
        N(x - a) & \leqslant \alpha ||x - a|| \\
        & \leqslant \alpha \eta ' \leqslant eta 
    \end{align*}

    Donc $N'(f(x) - \ell) \leqslant \gamma \varepsilon$, donc $||f(x) - \ell||' < \varepsilon$.
\end{proof}



\subsubsection{Cas où a appartient à A}

\begin{theoreme}{Continuité en un point}{}
    Soit $f : A \subset E \to E'$, et $a \in A$, $\ell \in E'$. 

    Si $f(x) \limitnorme{N'}{x}{a} \ell$. Alors $\ell = f(a)$.

    En ce cas on dit que $f$ est continue en $a$. Ainsi, $f$ continue en $a \in A$ ssi 
    $$\forall \varepsilon > 0, \, \exists \eta > 0, \, \forall x \in A, \, N(x - a) < \eta \Rightarrow N'(f(x) - f(a)) < \varepsilon$$
\end{theoreme}


\begin{proof}
    Supposons que $f(x) \limit x a \ell$. On a : 

    $\forall \varepsilon > 0, \, \exists \eta > 0, \, \forall x \in A, \, N(x - a) \leqslant \eta \Rightarrow N'(f(x) - \ell) \leqslant \varepsilon$.

    Soit $\varepsilon > 0$, et considérons $x = a$.

    On a alors bien $N(a - a) = 0 < \eta$, donc $N'(f(a) - \ell) \leqslant \varepsilon$.

    Or c'est vrai pour tout $\varepsilon > 0$, donc $N'(f(a) - \ell) = 0$, donc $f(a) = \ell$.
\end{proof}


\subsection{Extensions }

\begin{definition}{Extension notion de limite}{}
    On étend la notion de limite aux cas suivants : 
    \begin{enumerate}
        \item Si $E = \R, \, a = + \infty, \, (E', N')$ quelconque, $\ell \in E'$, et $A \subset \R$ 
            voisinage de $+ \infty$ et $f : A \to E'$.

            Alors $f(x) \limitnorme{N'}{x}{+ \infty} \ell$ ssi $\forall \varepsilon > 0, \, \exists c \in \R, \, \forall x \in A, \, x \geqslant c \Rightarrow N'(f(x) - \ell) \leqslant \varepsilon$.

        \item Si $E = \R, a = - \infty, \, \ell \in E'$, et $A \subset \R$ voisinage de $- \infty$, et $f : A \to E'$.

            Alors $f(x) \limitnorme{N'}{x}{- \infty} \ell$ ssi $\forall \varepsilon > 0, \, \exists c \in \R, \, \forall x \in A, \, x \leqslant c \Rightarrow N'(f(x) - \ell) \leqslant \varepsilon$.
        
        \item Si $E = E' = \R$, et $f(x) \limit x {\pm \infty} \pm \infty$ cf première année. 
        \item Soit $E$ un espace vectoriel normé, $A \subset E$, $a \in \overline A$, $E' = \R$ et $f: A \to \R$. Alors 
            $f(x) \limit x a + \infty$ ssi $\forall c \in \R, \, \exists \eta > 0, \, \forall x \in A, \, N(x - a) \leqslant \eta \Rightarrow f(x) \geqslant c$.

        \item Mêmes hypothèses. Alors $f(x) \limit x a - \infty$ ssi $\forall c \in \R, \, \exists \eta > 0, \, \forall x \in A, \, N(x - a) \leqslant \eta \Rightarrow f(x) \leqslant c$.
        \item Soit $A \subset E$ non borné, $f : E \to F, \, \ell \in F$. Alors $f(x) \limitnorme{N'}{x}{+ \infty} \ell$ ssi $\forall \varepsilon > 0, \, \exists c \in \R, \, \forall x \in A, \, N(x) \geqslant c \Rightarrow N'(f(x) - \ell) \leqslant \varepsilon$.
    \end{enumerate}
\end{definition}

\begin{theoreme}{Hors programme}{}
    Soit $f : A \subset E \to F$, $a \in \overline A$ (éventuellement $a = \pm \infty$ si $E = \R$), 
    $\ell \in E'$ (éventuellement $\ell = \pm \infty$ si $E' = \R$).

    Alors $f(x) \limitnorme{N'}{x}{a} \ell$ ssi l'image réciproque de tout voisinage de 
    $\ell$ est un voisinage relatif de $a$ vis à vis de $A$.
\end{theoreme}

\begin{remarque}
    Cette caractérisation concerne la définition et ses extensions, sauf la 6.
\end{remarque}

\idee{
    $\Rightarrow$ Utilisation de la limite avec des boules comme voisinages. 

    $\Leftarrow$ même idée dans l'autre sens. 
}

\begin{proof}
    Dans le cas $a \in E$, $\ell \in E'$. 

    Supposons que $f(x) \limitnorme{N'}{x}{a} \ell$.

    Soit $V$ un voisinage de $\ell$. Alors $\exists R > 0, \, B(\ell, R) \subset V$.

    Par définition de la limite avec $\varepsilon = \frac R 2$, $\exists \eta > 0, \, \forall x \in A, \, N(x - a) \leqslant \eta \Rightarrow N'(f(x) - \ell) \leqslant \frac R 2$.

    Donc si $x \in A \cap B(a, \eta)$, on a $f(x) \in V$. 

    Donc $A \cap B(a, \eta) \subset f^{-1}(V)$. 

    Donc $f^{-1}(V)$ est un voisinage relatif de $a$ vis à vis de $A$.

    \uindent{Réciproquement }

    \svspace Soit $\varepsilon > 0$. 

    Considérons $V = B'(\ell, \varepsilon)$. Donc $V$ est bien un voisinage de $\ell$. 

    Donc par hypothèse, $f^{-1}(V)$ est un voisinage relatif de $a$ vis à vis de $A$.

    Donc il existe $R > 0$ tel que $A \cap B(a, R) \subset f^{-1}(V)$.

    Posons $\eta = \frac R 2$. 

    Soit $x \in A$ tel que $d(x, a) \leqslant \eta$. On a $d(x, a) < R$, 
    donc $x \in A \cap B(a, R) $ donc $x \in f^{-1}(V)$, donc $f(x) \in V$, donc
    $d'(f(x), \ell) < \varepsilon$.
\end{proof}



\subsection{Caractérisation séquentielle }

\begin{theorement}{caractérisation séquentielle de la limite}{}
    Soit $A \subset E$ non vide, $a \in \overline A$, $\ell \in E'$. 

    Alors $f(x) \limitnorme{N'}{x}{a} \ell$ ssi pour toute suite $(u_n)_{n \in \N}$ 
    d'éléments de $A$ qui converge vers $a$, on a $f(u_n) \limitnorme{N'}{n}{\infty} \ell$.
\end{theorement}

\idee{
    $\Rightarrow$ Utilisation des defs et voisinages. 

    $\Leftarrow$ Raisonner par contraposée et absurde. 
}

\begin{proof}
    Supposons que $\forall \varepsilon > 0, \, \exists \eta > 0, \, \forall x \in A, \, N(x - a) \leqslant \eta \Rightarrow N'(f(x) - \ell) \leqslant \varepsilon$ (1).

    Soit $(u_n)_{n \in \N}$ une suite d'éléments de $A$ qui converge vers $a$.

    Alors $\forall \varepsilon > 0, \, \exists n_0 \in \N, \, \forall n \geqslant n_0, \, N(u_n - a) < \varepsilon$ (2)

    Soit $\varepsilon > 0$. Le (1) nous fournit un $\eta$, utilisons le dans le rôle de $\varepsilon$ dans (2),
    cela fournit alors un $n_0$.

    Soit $n \geqslant n_0$. 

    Alors on a (grâce au (2)) $N(u_n - a) \leqslant \eta$ donc (grâce au (1)) $N'(f(u_n) - \ell) \leqslant \varepsilon$.

    Donc $f(u_n) \limitnorme{N'}{n}{\infty} \ell$.

    \uindent{Réciproquement } raisonnons par contraposée. 

    \svspace Supposons que $\exists \varepsilon > 0, \, \forall \eta > 0, \, \exists x \in A, \, d(x, a) \leqslant \eta $ et $d'(f(x), \ell) > \varepsilon$ (3).

    Soit $n \in \N^*$. Cette phrase avec $\frac 1 n$ dans le rôle de $\eta$ fournit
    un $x_n \in A$ tel que $N(x_n - a) \leqslant \frac 1 n$ et $N'(f(x_n) - \ell) > \varepsilon$.

    On a $N(x_n - a) \limit n \infty 0$, donc $x_n \limitnorme N n \infty a$.

    On a de plus $N'(f(x_n) - \ell) > \varepsilon$. Si on avait $f(x_n) \limitnorme{N'}{n}{\infty} \ell$, 
    on aurait $\varepsilon$ la limite, donc $0 \geqslant \varepsilon$, absurde.

    Donc $f(x_n)$ ne converge pas vers $\ell$.
\end{proof}


\begin{theoreme}{Carac séquentielle continuité}{}
    Soit $A \subset E$ non vide, $a \in A$. 

    Alors $f$ est continue en $a$ ssi pour toute suite $(u_n)_{n \in \N}$ d'éléments de $A$ 
    qui converge vers $a$ on a $f(u_n) \limitnorme{N'}{n}{\infty} f(a)$.
\end{theoreme}

\subsection{Opérations}

\begin{theoreme}{}{}
    Soient $f, g : A \subset E \to F$, $\lambda : A \to \K$, $\alpha \in K$ et $\ell_1, \ell_2 \in F$. 

    Si $f(x) \limitnorme{N'}{x}{a} \ell_1$ et $g(x) \limitnorme{N'}{x}{a} \ell_2$, alors
    $(f + g)(x) \limit x a \ell_1 + \ell_2$ et $\alpha f(x) \limit x a \alpha \ell_1$

    Si $f(x) \limit x a \ell_1$ et $\lambda(x) \limit x a \alpha$, alors $(\lambda f)(x) \limit x a \alpha \ell_1$
\end{theoreme}


\begin{proof}
    Par caractérisaiton séquentielle. Soit $u_n \longrightarrow a$. 

    Supposons que $f(x) \limit x a \ell_1$ et $g(x) \limit x a \ell_2$ alors 
    $f(u_n) \limit n {+ \infty} \ell_1$ et $g(u_n) \longrightarrow \ell_2$. 

    Par caractérisation séquentielle. Donc $(f + g)(u_n) \longrightarrow \ell_1 + \ell_2$. 

    Donc (carac seq), $(f + g)(x) \longrightarrow \ell_1 + \ell_2$

    De même pour les autres. 
\end{proof}

Conséquence : soient $f, g : A \subset E \to F$, $\lambda : A \to \K$, $\alpha \in \K, a \in A$. 

Si $f$ et $g$ dont continues en $a$, $f + g$ et $\alpha f$ aussi. 

Si $f$ et $\lambda$ continues en a, $\lambda f$ aussi. 

\begin{proposition}{}{}
    Soit $A \subset E$, $a \in \overline A$, $B \subset E'$, $b \in \overline B$ et 
    $(E", N")$ un evn, et soit $\ell \in E"$, $f : A \to E'$ tel que $\forall x \in A, f(x) \in B$, $g: B \to E"$. 

    Si $f(x) \limitnorme{N'}x a b \text{ et } g(y) \limitnorme{N"}y a \ell$ alors 

    $$(g \circ f )(x) \limitnorme{N"} x a \ell$$
\end{proposition}

\begin{proof}
    $\forall \varepsilon > 0, \exists \eta_1 > 0, \forall x \in A, N(x - a) \leqslant \eta_1 \Rightarrow N'(f(x) - b) \leqslant \varepsilon$ (1)

    $\forall \varepsilon > 0, \exists \eta_2 > 0, \forall y \in B, N(y - b) \leqslant \eta_2 \Rightarrow N"(g(y) - \ell) \leqslant \varepsilon$

    Soit $\varepsilon > 0$, (2) fournit $\eta$. Utilisons le (1) avec $\eta_2$ dans le rôle 
    de $\varepsilon$. Cela fournit $\eta_1$

    Soit $x \in A$ te qie $N(x - a) \leqslant \eta_1$ alors $N'(f(x) - b) \leqslant \eta_2$ par 1

    Donc $N"(g(f(x)) - \ell) \leqslant \varepsilon$ par 2 avec $y = f(x)$
\end{proof}

Conséquence : 

Soit $f : A \subset E \to E'$ tel que $f(A) \subset B$, $g : B \to E"$ $a \in A$. Si $f$ continue en $a$ et $g$ continue en $f(a)$ alors $g \circ f$ continue en a 


\subsection{Cas des espaces vectoriels normés produits}

Soient $(E'_1, N'_1), \ldots, (E'_p, N'_p)$ des evn, et $E' = E'_1 \times \cdots \times E'_p$.
muni de la norme produit $N'_\infty$, et $\fonction{f}{A\subset E}{E'}{x}{f(x) = (f_1(x), \ldots, f_p(x))}$.


\begin{proposition}{Limite espace produit}{}
    Soit $a \in \overline A$, $\ell \in E', \, \ell = (\ell_1, \ldots, \ell_p)$.

    Alors $f(x) \limitnorme{N'_\infty}{x}{a} \ell$ ssi $\forall j \in \ninter{1}{p}, \, f_j(x) \limitnorme{N'_j}{x}{a} \ell_j$.

\end{proposition}

\uindent{Conséquence }

Soit $a \in A$. Alors $f$ est continue en $a$ ssi $\forall j \in \ninter{1}{p}, \, f_j$ est continue en $a$.

\idee{
    Caractérisation séquentielle de la limite.
}

\begin{proof}
    $f(x) \limitnorme{N'_\infty}{x}{a} \ell$ ssi pour toute suite $(u_n)_{n \in \N}$ 
    d'éléments de $A$ qui converge vers $a$, on a $f(u_n) \limitnorme{N'_\infty}{n}{\infty} \ell$.
    ssi pour toute suite $(u_n)_{n \in \N}$ d'éléments de $A$ qui converge vers $a$, 
    $\forall j \in \ninter{1}{p}, \, f_j(u_n) \limitnorme{N'_j}{n}{\infty} \ell_j$.

    ssi $\forall j \in \ninter{1}{p}, \, f_j(x) \limitnorme{N'_j}{x}{a} \ell_j$.
\end{proof}



\section{Continuité sur une partie}

\begin{definition}{Continuité restriction sur une partie }{}
    Soit $f: A \to E$, et soit $B \subset A$. On dit que $f$ est continue sur $B$ ssi 
    $\forall a \in B, \, f_{|B}$ est continue en $a$.
\end{definition}


\begin{definition}{Continuité sur un ensemble}{}
    Soit $f : A \to E$. On dit en particulier que $f$ est continue sur $A$ ssi
    $\forall a \in A, \, f$ est continue en $a$.
\end{definition}


\begin{proposition}{Continuité sur un ensemble}{}
    Soit $f : E \to E'$ et $U \subset E$ un ouvert.

    Alors $f$ est continue sur $U$ ssi $\forall a \in A$, $f$ est continue en $a$. 
\end{proposition}

\begin{proof}
    $f$ est continue sur $U$ ssi $\forall a \in U, \, f_{|U}$ est continue en $a$ 
    
    ssi $\forall a \in U, \, \forall (u_n)_{n \in \N}$ suite d'éléments de $U$ qui converge vers $a$,
    $f_{|U}(u_n) \limitnorme{N'}{n}{\infty} f(a)$

    ssi $\forall a \in U, \, \forall (u_n)_{n \in \N}$ suite d'éléments de $E$ qui converge vers $a$,
    $f(u_n) \limitnorme{N'}{n}{\infty} f(a)$ car comme $U$ est ouvert, pour $n$ assez
    grand, si $u_n \in E$ et $u_n \to a$, alors $u_n \in U$.

    ssi $\forall a \in U, \, f$ est continue en $a$.
\end{proof}

\subsection{Théorèmes d'opérations}

\begin{theorement}{vfv}{}
    Soit $f, g, : A \to E$, $\lambda : A \to \K$, $\alpha \in \K$.

    Si $f$ et $g$ sont continues sur $A$, alors $\alpha f + g$ et $\lambda f$ sont 
    continues sur $A$. 
\end{theorement}


\uindent{Conséquence }

L'ensemble $\mathscr C(A, E')$ des applications continues de $A$ dans $E'$ est un $\K$-espace vectoriel.

\begin{theorement}{ff}{}
    Soit $f : A \to E'$, $g : B \subset E' \to E"$ tel que $f(A) \subset B$.

    Si $f \in \mathscr C(A, E')$ et $g \in \mathscr C(B, E")$, alors 
    $g \circ f \in \mathscr C(A, E")$.
\end{theorement}

\begin{theoreme}{Égalité d'applications avec la densité}{}
    Deux applications continues qui coincident sur une partie dense sont égales. 
\end{theoreme}

\idee{
    Caractérisation séquentielle, on utilise une suite par point de $A$. 
}

\begin{proof}
    Soient $f, g \in \mathscr C(A, E')$ et $B \subset A$ dense dans $A$ telles que
    $A \subset \overline B$ telle que $\forall x \in B, \, f(x) = g(x)$.

    Soit $x \in A$. 

    $x \in \overline B$, donc il existe une suite $(u_n)_{n \in \N}$ d'éléments de $B$ qui converge vers $x$.
    
    Comme $u_n \in B$, on a $\forall n \in \N, \, f(u_n) = g(u_n)$ et alors par 
    passage à la limite, $f(x) = g(x)$ car $f$ et $g$ sont continues en $x$.

    Donc $f = g$
\end{proof}

\begin{exemple}
    Soit $f \in \mathscr C(\R, \R)$ telle que $\forall x, y \in \R, \, f(x + y) = f(x) + f(y)$.

    Montrons que $\forall x \in \R, \, f(x) = f(1) x$.

    Posons $\forall x \in \R, \, g(x) = f(1) x$. Donc $g \in \mathscr C(\R, \R)$ et 
    par récurrence, on a $\forall n \in \N, \, f(n) = g(n)$.

    Soit $p \in \Z$ tel que $p < 0$. Alors $f(-p) + f(p) = f(p - p) = f(0)$

    Donc $f(-p) = - f(p)$, donc $\forall p \in \Z, \, f(p) = g(p)$.

    Donc égalité pour $x \in \Z$. 

    Soit $x \in \Q$. Il existe $p \in \Z, \, q \in \N^*$ tels que $x = \frac p q$.

    Alors : 
    \begin{align*}
        qf(x) & = f(qx) \\
        & = f(p) \\
        & = g(p) = pf(1) = q \frac p q f(1) \\
        & = q g(x)
    \end{align*}

    Donc $f(x) = g(x)$ pour $x \in \Q$.

    Or $\Q$ est dense dans $\R$, donc par le théorème, $f = g$.
\end{exemple}

\subsection{Continuité et image réciproque }

\begin{theoreme}{Nature image réciproque fonction continue }{}
    Soit $f \in \mathscr C(A, E')$. Alors l'image réciproque d'un ouvert (resp. fermé) 
    de $E'$ par $f$ est un ouvert (resp. fermé) relatif de $A$.

    Soit $f \in \mathscr C(E, E')$. Alors l'image réciproque d'un ouvert (resp. fermé)
    de $E'$ par $f$ est un ouvert (resp. fermé) de $E$.
\end{theoreme}

\begin{remarque}
    La réciproque est vraie mais pas au programme.
\end{remarque}

\begin{proof}
    Soit $f \in \mathscr C(A, E')$, soit $B \subset E'$ fermé. 

    Montrons que $f^{-1}(B)$ est un fermé relatif de $A$.

    Considérons $(u_n) _{n \in \N}$ une suite d'éléments de $f^{-1}(B)$ qui converge vers $\ell \in A$.

    $f$ est continues en $\ell$, donc $f(u_n) \rightarrow f(\ell)$. 
    Or $\forall n, \, f(u_n) \in B$ et $B$ est fermé, donc $f(\ell) \in B$.

    Donc $\ell \in f^{-1}(B)$. 

    Donc $f^{-1}(B)$ est un fermé relatif de $A$.

    \uindent{Pour un ouvert }

    \svspace Soit $U$ un ouvert de $E'$.

    Prenons $B = E' \setminus U$, alors $f^{-1}(U) \subset A$. 

    Donc $f^{-1}(U) = A \setminus f^{-1}(B)$ or $f^{-1}(B)$ est un fermé relatif de $A$.

    Donc $f^{-1}(U)$ est un ouvert relatif de $A$.
\end{proof}

\subsection{Applications uniformément continues et lipschitziennes}

\begin{definition}{Uniforme continuité}{}
    Soit $f : A \subset E \to E'$. On dit que $f$ est uniformément continue sur $A$ ssi 
    $$\forall \varepsilon > 0, \, \exists \eta > 0, \, \forall x, y \in A, \, N(x - y) \leqslant \eta \Rightarrow N'(f(x) - f(y)) \leqslant \varepsilon$$
\end{definition}

\begin{definition}{Lipschitzienne}{}
    Soit $f : A \subset E \to E'$ et $k \geqslant 0$. On dit que $f$ est lipschitzienne 
    de rapport $k$ ssi 
    $$\forall x, y \in A, \, N'(f(x) - f(y)) \leqslant k N(x - y)$$

    On dit que $f$ est lipschitzienne ssi elle est lipschitzienne de rapport $k$ pour un certain $k \geqslant 0$.
\end{definition}


\begin{remarque}
    La caractère lipschitzien est invariant si on change les normes en des normes 
    équivalentes, mais le rapport peut changer.
\end{remarque}

\begin{proposition}{Lien entre lipschitzienne, uniforme continuité et continuité}{}
    Toute applications lipschitzienne est uniformément continue. 

    Toute application uniformément continue est continue.
\end{proposition}

\idee{
    Poser le bon $\eta$.
}

\begin{proof}
    Soit $f$ $k$ lipschitzienne de $A$ dans $E'$.

    Soit $\varepsilon > 0$, posons $\eta = \frac \varepsilon k$.

    Soient $x, y \in A$ tels que $N(x - y) \leqslant \eta$.

    alors $N'(f(x) - f(y)) \leqslant k N(x - y) \leqslant k \eta = \varepsilon$.

\end{proof}

\begin{proposition}{Nature applications coordonnées}{}
    Dans un evn de dimension finie, les applications coordonnées sont lipschitziennes, 
    donc uniformément continues, donc continues. 
\end{proposition}


\begin{proof}
    $E$ evn de base $B = (e_1, \ldots, e_n)$. Munissons le de la norme $||\cdot||_\infty$.

    Soit $i \in \ninter{1}{n}$, et soit $\fonction{c_i}{E}{\K}{x = \sum_{j=1}^n x_j e_j}{x_i}$.

    Soient $x, y \in E$. Alors $|c_i(x) - c_i(y)| = |c_i(x - y) | \leqslant ||x - y||_\infty$.

    Donc $c_i$ est 1-lipschitzienne. 
\end{proof}


\begin{theoreme}{Lien dérivée et lipschitzienne}{}
    Soit $f \in \mathscr C^1(I, \R)$, alors $f$ est lipschitzienne ssi $f'$ est bornée.
\end{theoreme}


\begin{proof}
    Supposons que $f$ est $k$-lipschitzienne. 

    Soit $x \in I$ tel que $\forall y \in I$ avec $x \neq y$, $\left|\frac{f(x) - f(y) }{x - y}\right| \leqslant k$. 

    Et en faisant tendre $y$ vers $x$, on obtient $|f'(x)| \leqslant k$.
\end{proof}

\begin{proposition}{Nature fonction distance à une partie}{}
    Soit $A \subset E$ non vide, et $\fonction{f}{E}{R}{x}{d_A(x)}$. 

    Alors $d_A$ est 1-lipschitzienne. 
\end{proposition}

\begin{proof}
    Soient $x, y \in E$. 

    Notons $D_x = \{d(x, a), a \in A\}$ et $D_y = \{d(y, a), a \in A\}$.

    Soit $a \in A$. 

    On a $N(x - a) \leqslant N(x - y) + N(y - a)$, donc $N(x - a) \in D_x$ et 
    $d_A(x)$ en est un minorant donc $\forall a \in A, \, d_A(x) \leqslant N(x - y) + \underbrace{N(y - a)}_{\in D_y}$.

    Et donc $\underbrace{d_A(x) - N(x - y)}_{\text{minorant de } D_y} \leqslant N(y - a)$

    Donc $d_A(x) - N(x - y) \leqslant d_A(y)$.

    Donc $d_A(x) - d_A(y) \leqslant N(x - y)$, et par symétrie des rôles, 

    $|d_A(x) - d_A(y)| \leqslant N(x - y)$.
\end{proof}


\section{Continuité des applications (multi)linéaires}
\subsection{Théorème : critère de continuité }

\begin{theorement}{critère de continuité pour une application linéaire}{}
    Soient $(E, N)$ et $(E', N')$ deux evn, et $ f \in \mathscr L(E, E')$.

    Alors $f$ est continue sur $E$ ssi 
    $$\exists c > 0, \, \forall x \in E, \, N'(f(x)) \leqslant c N(x)$$
\end{theorement}

\idee{
    $\Rightarrow$ Utilisation de la continuité en 0 et du fait que $f$ est linéaire.

    $\Leftarrow$ Montrer $f$ lipschitzienne $\eta$.
}

\begin{proof} 
    $\Rightarrow$ Supposons que $f$ est continue sur $E$. En particulier, $f$ est 
    continue en $0_E$. Considérons $\varepsilon = 1$.

    Cela fournit $\eta > 0$ tel que $\forall x \in E, \, N(x - 0) \leqslant \eta \Rightarrow N'(f(x) - f(0)) \leqslant 1$.

    Or $f$ linéaire, donc $f(0_E) = 0_{E'}$ donc 
    
    $$N(x) \leqslant \eta \Rightarrow N'(f(x)) \leqslant 1$$

    Soit $x \in E \setminus \{0\}$. 
    
    Posons $\displaystyle y = \eta \frac {x}{N(x)} $.

    Alors $N(y) = \eta$, donc $N'(f(y)) \leqslant 1$.

    Donc $N'\left( \frac{\eta }{N(x)} f(x) \right) \leqslant 1$ car $f$ linéaire, 
    
    et par homogénéité de $N'$, $ \frac{\eta }{N(x)} N'(f(x)) \leqslant 1$.

    Donc $N'(f(x)) \leqslant \frac{N(x)}{\eta}$.

    Comme $N'(f(0_E)) = N(0_{E'}) = 0 \leqslant \frac 1 \eta 0 $, 
    
    en posant $c = \frac 1 \eta$, on a bien le résultat. 

    \lvspace $\Leftarrow$ Supposons l'inégalité et montrons que $f$ est $c$-lipschitzienne. 

    Soient $x, y \in E$. Alors $N'(f(x) - f(y)) = N'(f(x - y)) $ car $f$ linéaire.

    Donc $N'(f(x) - f(y)) \leqslant c N(x - y)$ par hypothèse. 

    Donc $f$ est $c$-lipschitzienne, donc continue sur $E$.
\end{proof}

\begin{remarque}
    On a montré au passage que $f$ est continue sur $E$ ssi $f$ est continue en $0$ ssi $f$ est lipschitzienne.  

    Pour montrer que $f$ n'est pas continue, on exhibe une suite $(x_n)_{n \in \N}$ de $E$ 
    telle que $\displaystyle \frac{N'(f(x_n))}{N(x_n)} \to + \infty$. Alors $\forall c > 0, \, \exists x_n \in E, \, N'(f(x_n)) > c N(x_n)$, 
    donc $f$ non continue.
\end{remarque}

\begin{exemple}
    Soit $\fonction f {\R[X]} \R P {P(1)}$, $f$ est linéaire. 

    \avspace En munissant $\R[X]$ de la norme $||P||_\infty = \max\{P(t), t \in [0, 1]\}$, on a 
    $\forall P \in \R[X], \, |f(P)| = |P(1)| \leqslant ||P||_\infty$.

    Donc par critère de continuité des applications linéaires, $f$ est continue.

    \avspace En munissant $\R[X]$ de la norme $N_\infty = \max\{P(t), t \in \inter 0 {\frac 1 2}\}$

    Posons $P_n = X^n$. Alors $| f(P_n) | = 1$ et $||P_n||_\infty = \max\{P_n(t), t \in \inter 0 {\frac 1 2}\} = \left(\frac 1 2\right)^n$.
\end{exemple}

\begin{definition}{Applications linéaires continues}{}
    On pose $\mathscr L_c(E, E') = \mathscr L (E, E') \cap \mathscr C(E, E')$ l'espace vectoriel 
    des applications linéaires continues de $E$ dans $E'$.
\end{definition}

\begin{theoreme}{Continuité des applications linéaires en dimension finie}{}
    Si $E$ est de dimension finie, alors $\mathscr L_c(E, E') = \mathscr L(E, E')$.

    Ainsi toute application linéaire au départ d'un espace de dimension finie est continue
\end{theoreme}

\begin{proof}
    Soit $E$ un $\K$-espace vectoriel de dimension finie et de base 
    $B = (e_1, \ldots, e_p)$. Munissons $E$ de la norme $||\cdot||_\infty$ asociée à $B$.

    Soit $(E', N')$ un $\K$ evn. 

    Soit $f \in \mathscr L(E, E')$, et soit $x \in E$. 

    Décomposons le en $\displaystyle x = \sum_{i=1}^p x_i e_i$. Alors : 

    \begin{align*}
        N'(f(x)) & = N'\left( f\left( \sum_{i=1}^p x_i e_i \right) \right) \\
        & = N'\left( \sum_{i=1}^p x_i f(e_i) \right) \\
        & \leqslant \sum_{i=1}^p N(x_i f(e_i)) \\
        & \leqslant \sum_{i=1}^p |x_i| N'(f(e_i)) \\
        & \leqslant \underbrace{\left( \sum_{i=1}^p N'(f(e_i)) \right)}_{ = K \geqslant 0} ||x||_\infty
    \end{align*}

    Donc par critère de continuité des applications linéaires, $f$ est continue.
\end{proof}

\uindent{Conséquence } 

En dimension finie, les applications coordonnées sont continues car linéaires, donc 
tout polynôme en les coordonnées est continu. 


\avspace $\fonction{}{\mathscr M_n(\K)} \K A {\det(A)}$ est continue.

$\fonction {}{\mathscr M_{n, p}(\K) \times \mathscr M_{p, q}(\K)}{\mathscr M_{n, q}(\K)}{(A, B)}{AB}$ est continue.

\begin{proof}
    Le déterminant est polynôme en les coordonées de la matrice dans la base cononique des 
    matrices élémentaires. 

    \svspace Chacune des coordonnées de $AB$ est un polynôme en les coordonnées de $A$ et $B$.
\end{proof}


\subsection{Norme subordonnée d'une application linéaire continue }

\begin{definition}{Norme d'opérateur}{}
    Soient $(E, N)$, $(E', N')$ deux evn sur $\K$ et $f \in \mathscr L_c(E, E')$.

    Alors par critère de continuité, $\left\{ \frac{N'(f(x))}{N(x)}, x \in E \setminus \{0\} \right\}$ est
    une partie de $\R^+$ majorée. 

    On appelle norme d'opération (ou norme subordonnée, ou parfois norme triple) de $f$ le réel : 
    $$||f||_{op} = |||f||| = \sup \left\{ \frac{N'(f(x))}{N(x)}, x \in E \setminus \{0\} \right\} $$ 


    On a en particulier $\forall x \in E, \, N'(f(x)) \leqslant |||f||| \cdot N(x)$.
\end{definition}


\begin{exemple}
    $\fonction u E \R \varphi {\int_0^1 \varphi}$ donc $u \in \mathscr L(E, \R)$. 

    Donc $|u(\varphi)| = \left| \int_0^1 \varphi(t) dt \right| \leqslant \int_0^1 |\varphi(t)| dt \leqslant ||\varphi||_\infty$.

    Donc $u$ est continue. 

    Posons $\varphi = 1$, alors $|u(\varphi)| = 1 = ||\varphi||_\infty$.

    Donc le majorant 1 de $\left\{ \frac{|u(\varphi)|}{||\varphi||_\infty}, \varphi \in E \setminus \{0\} \right\}$ est atteint. 

    Donc $|||u||| = 1$.
\end{exemple}


\begin{proposition}{}{}
    Soit $f \in \mathscr L_c(E, E')$ alors $|||f||| = \sup \{ N'(f(x)), \, x \in S(0, 1) \}$
\end{proposition}

\begin{proof}
    Notons $A = \{\frac{N'(f(x))}{N(x)}, \, x \in E \setminus \{0\}\}$,
    $B = \{N'(f(x)), \, x \in E \mid N(x) = 1\}$.

    Soit $\alpha \in A$, il existe $x \in E $ tel que $x \neq 0$ et $\alpha = \frac{N'(f(x))}{N(x)}$.

    Posons $y = \frac x {N(x)}$, alors : 

    \begin{align*}
        N'(f(y)) & = N'\left( f\left( \frac x {N(x)} \right) \right) \\
        & = N'\left( \frac{1}{N(x)} f(x) \right) \\
        & = \frac{1}{N(x)} N'(f(x)) \\
        & = \alpha
    \end{align*}

    Donc $\alpha \in B$. 

    Il existe $y \in S(0, 1)$, $\beta = N'(f(y))$. 

    Posons $x = y$, on a $\beta = \frac{N'(f(y))}{N(y)}$. 

    Donc $\beta \in A$.
\end{proof}


\begin{theoreme}{Norme d'opérateur est une norme }{}
    La norme d'opérateur $|||\cdot|||$ est une norme sur $\mathscr L_c(E, E')$.
\end{theoreme}


\begin{proof}
    \begin{itemize}
        \item $|||\cdot|||$ est positive par définition.
        \item Supposons que $|||f||| = 0$, alors $\forall x \in E \setminus \{0\}, \, N'(f(x)) = 0$.
            Donc $f(x) = 0_{E'}$ car $N'$ est une norme. Donc $f = 0_{\mathscr L(E, E')}$ car linéaire. 
        \item Soit $\lambda \in \K$
    \end{itemize}

    $|||\lambda f||| = \sup \left\{ \frac{N'(\lambda f(x))}{N(x)}, x \in E \setminus \{0\} \right\} = \sup \left\{ |\lambda| \frac{N'(f(x))}{N(x)}, x \in E \setminus \{0\} \right\} = |\lambda| |||f|||$.

    \begin{itemize}
        \item Soient $f, g \in \mathscr L_c(E, F)$.
    \end{itemize}

    Soit $x \in E$ tel que $x \neq 0$.

    \begin{align*}
         \frac{N'((f + g)(x))}{N(x)} & = \frac{N'(f(x) + g(x))}{N(x)} \\
            & \leqslant \frac{N'(f(x)) }{N(x)} + \frac{N'(g(x))}{N(x)} \\
            & \leqslant |||f||| + |||g|||
    \end{align*}

    Donc par def du sup, $|||f + g||| \leqslant |||f||| + |||g|||$.
\end{proof}


\begin{proposition}{Sous-multiplicativité de la norme d'opérateur }{}
    Soient $(E, N)$, $(E', N')$, $(E", N")$ trois evn sur $\K$, et soient $f \in \mathscr L_c(E, E')$ et $g \in \mathscr L_c(E', E")$.

    Alors $g \circ f \in \mathscr L_c(E, E")$ et $|||g \circ f||| \leqslant |||g||| \cdot |||f|||$.
\end{proposition}

\begin{proof}
    Pour $x \in E, \, N'(f(x)) \leqslant |||f||| N(x)$ vrai si $x \neq 0$ par def 
    de la norme d'opérateur, et si $x = 0$, $N'(f(0)) = N'(0) = 0$.
    Donc  

    \begin{align*}
        N"(g(f(x))) & \leqslant |||g||| N'(f(x)) \\
        & \leqslant |||g||| \cdot |||f||| N(x)
    \end{align*}

    Donc $|||g||| \cdot |||f|||$ est un majorant de $\left\{ \frac{N"(g\circ f(x))}{N(x)}, x \in E \setminus \{0\} \right\}$.

    Donc $|||g \circ f||| \leqslant |||g||| \cdot |||f|||$.
\end{proof}

\subsubsection{Cas des endomorphismes }

\begin{proposition}{Norme d'algèbre}{}
    Pour $f \in \mathscr L_c(E)$, on a $||||f||| = \sup \left\{ \frac{N(f(x))}{N(x)}, \, x \in E \setminus \{0\} \right\}$
    et $|||f|||$ est une norme qui vérifie : 

    $$\forall f, g \in \mathscr L_c(E), \, |||f \circ g||| \leqslant |||f||| \cdot |||g|||$$

    On dit que $|||\cdot|||$ est une norme d'algèbre. 
\end{proposition}


\subsubsection{Cas des matrices }

\begin{proposition}{Norme matricielle}{}
    On se place dans $\mathscr M_n(\K)$ et on munit $\K^n$ d'une norme $N$.

    On définit pour $A \in \mathscr M_n(\K)$, $|||A||| = \sup \left\{ \frac{N(AX)}{N(X)}, X \in \K^n \setminus \{0\} \right\}$
    
    Alors $|||\cdot|||$ est une norme d'algèbre sur $\mathscr M_n(\K)$, en particulier 
    $$\forall A, B \in \mathscr M_n(\K), \, |||AB||| \leqslant |||A||| \cdot |||B|||$$

    On dit aussi que $|||\cdot|||$ est une norme matricielle. 
\end{proposition}

\begin{remarque}
    Une norme équivalente à une norme matricielle n'est pas forcément une norme matricielle.
\end{remarque}

\begin{proof}
    On identifie $\K^n = \mathscr M_{n, 1}(\K)$ matrices colonnes qui existe 
    car $\fonction{}{\K^n}{\K^n}{X}{AX}$ est linéaire et continue car $\K^n$ est de dimension finie. 
    
\end{proof}



\begin{exemple}
    Posons $N_\infty(A) = \max\{ |a_{i, j}|, \, i, j \in \ninter{1}{n}\}$. 
    C'est bien une norme, mais elle n'est pas matricielle. 

    En effet, posons 
    $A = \begin{pmatrix}
        1 & \dots & 1 \\
        \vdots & \ddots & \vdots \\
        1 & \dots & 1
    \end{pmatrix}$ alors $N_\infty(A) = 1$.

    On a $A^2 = nA$, donc $N_\infty(A^2) = n$ et donc $||A^2||_\infty = n > 1 = ||A||_\infty^2$. 

    \lvspace Posons $N(A) = n \cdot N_\infty(A)$, qui est bien une norme, équivalente à $N_\infty$.

    Montrons que $N$ est une norme matricielle.

    Soient $A, B \in \mathscr M_n(\K)$, et $C = AB$. 
    \begin{align*}
        |c_{i, j}| & = \left| \sum_{k=1}^n a_{i, k} b_{k, j} \right| \\
        & \leqslant \sum_{k=1}^n |a_{i, k}| |b_{k, j}| \\
        & \leqslant n N_\infty(A) N_\infty(B)
    \end{align*}

    Donc 
    \begin{align*}
        N(C) & = n N_\infty(C) \\
        & \leqslant n^2 N_\infty(A) N_\infty(B) \\
        &  \leqslant N(A) N(B)
    \end{align*}


\end{exemple}


\subsubsection{LHP, étude de $||| \cdot |||_1$ associée à $|| \cdot ||_1$}


Soit $A \in \mathscr M_n(\K)$, $|||A|||_1 = \sup (\{ ||AX||_1, X \in \K^n, ||X||_1 = 1 \} = z) $.

Soit $X \in \K^n$ tel que $||X||_1 = 1$, c'est à dire $X = (x_1, \ldots, x_n)^T$ et $\sum_{i=1}^n |x_i| = 1$.

$AX = \begin{pmatrix}
    y_1 \\
    \vdots \\
    y_n
\end{pmatrix}$  avec $||AX||_1 = \sum_{i=1}^n |y_i|$, 
or $\displaystyle \forall i, \, y_i = \sum_{j=1}^n a_{i, j} x_j$. 

Donc 
\begin{align*}
    || AX ||_1 & = \sum_{i=1}^n \left| \sum_{j=1}^n a_{i, j} x_j \right| \\
    & \leqslant \sum_{i=1}^n \sum_{j=1}^n |a_{i, j}| |x_j| \\
    & \leqslant \sum_{j=1}^n |x_j| \sum_{i=1}^n |a_{i, j}| 
\end{align*}


Posons $k = \max \left\{ \sum_{i=1}^n |a_{i, j}|, j \in \ninter{1}{n} \right\}$.

Alors $||AX||_1 \leqslant k \sum_{j=1}^n |x_j| = k ||X||_1 = k$.


Soit $j_0$ un indice tel que $k = \sum_{i=1}^n |a_{i, j_0}|$, et posons pour 
$j \neq j_0, \, x_j = 0$ et $x_{j_0} = 1$ on a bien $||X||_1 = 1$ et
$||AX||_1 = \sum_{i=1}^n |a_{i, j_0} x_{j_0}| = k$.

Donc $k$ est un majorant de $z$ qui est atteint en $X$,

Donc $k = |||A|||_1$.

Donc $|||A|||_1 = \max \left\{ \sum_{i=1}^n |a_{i, j}|, j \in \ninter{1}{n} \right\}$.


\subsubsection{LHP, étude de $||| \cdot |||_\infty$ associée à $|| \cdot ||_\infty$}

Soit $A \in \mathscr M_n(\K)$, $|||A|||_\infty = \sup (\{ ||AX||_\infty, X \in \K^n, ||X||_\infty = 1 \} = z) $.

Soit $X \in \K^n$ tel que $||X||_\infty = 1$, c'est à dire $X = (x_1, \ldots, x_n)^T$ et $\max(|x_1|, \ldots, |x_n|) = 1$.

$AX = \begin{pmatrix}
    y_1 \\
    \vdots \\
    y_n
\end{pmatrix}$  avec $||AX||_\infty = \max(|y_1|, \ldots, |y_n|)$, 
or $\displaystyle \forall i, \, y_i = \sum_{j=1}^n a_{i, j} x_j$. Donc 

\begin{align*}
    |y_i| & = \left| \sum_{j=1}^n a_{i, j} x_j \right| \\
    & \leqslant \sum_{j=1}^n |a_{i, j}| |x_j| \\
    & \leqslant \sum_{j=1}^n |a_{i, j}| ||X||_\infty \leqslant \sum_{j=1}^n |a_{i, j}|
\end{align*}

Donc $\displaystyle ||AX||_\infty = \max \left\{ \sum_{j = 1 }^{n } |a_{i, j}|, i \in \ninter 1 n  \right\} = k$

Soit $i_0$ un indice tel que $k = \sum_{j=1}^n |a_{i_0, j}|$, et posons 
$X = (x_1, \ldots, x_n)^T$. Notons $\forall j, \, a_{i_0, j} = |a_{i_0, j}| e^{i \theta_j}$ 

Posons $x_j = e^{-i \theta_j}$, alors $||X||_\infty = 1$.

Alors $|y_{i_0}| = \left|\sum_{j = 1 }^{n } | a_{i_0, j} | e^{i \theta_j e^{-i \theta_j}}\right| = \sum_{j=1}^n |a_{i_0, j}| = k$.

Donc $|| AX||_\infty \geqslant k$, or $k$ majorant donc $k = |||AX |||_\infty$ et $k = \sup Z$. 


\subsection{Applications multilinéaires }

\begin{theoreme}{Continuité application $p$-linéaire}{}
    Soient $(E_1, N_1), \ldots, (E_p, N_p)$ et $(E', N')$ des evn et 
    
    $\fonction{f}{E_1 \times \cdots \times E_p}{E'}{(x_1, \ldots, x_p)}{f(x_1, \ldots, x_p)}$ 
    $p$-linéaire. 

    Alors $f$ est continue ssi $$\exists c > 0, \, \forall (x_1, \ldots, x_p) \in E_1 \times \cdots \times E_p, \, N'(f(x_1, \ldots, x_p)) \leqslant c \prod_{i=1}^p N_i(x_i)$$
\end{theoreme}


Démo non exigible 


\begin{proposition}{Continuité application $p$-linéaire en dimension finie}{}
    Soient $(E_1, N_1), \ldots, (E_p, N_p)$ des evn de dimension finie et $(E', N')$ un evn.

    Alors si $f : E_1 \times \cdots \times E_p \to E'$ est $p$-linéaire, alors $f$ est continue.
\end{proposition}

\begin{proof}
    Soit $B_1 = (e_{1, 1}, \ldots, e_{1, n_1})$ une base de $E_1$, ..., $B_p = (e_{p, 1}, \ldots, e_{p, n_p})$ une base de $E_p$.

    Soient $\displaystyle x_1 = \sum_{i=1}^{n_1} a_{1, i_1} e_{1, i_1} \in E_1$, ..., $\displaystyle x_p = \sum_{i=1}^{n_p} a_{p, i_p} e_{p, i_p} \in E_p$.

    Alors 
    
    \begin{align*}
        f(x_1, \ldots, x_p) & = f\left( \sum_{i_1=1}^{n_1} a_{1, i_1} e_{1, i_1}, \ldots, \sum_{i_p=1}^{n_p} a_{p, i_p} e_{p, i_p} \right) \\
        & = \sum_{i_1=1}^{n_1} \cdots \sum_{i_p=1}^{n_p} a_{1, i_1} \cdots a_{p, i_p} f(e_{1, i_1}, \ldots, e_{p, i_p}) 
    \end{align*}

    Donc 
    \begin{align*}
        N'(f(x_1, \ldots, x_p)) & \leqslant \sum_{i_1=1}^{n_1} \cdots \sum_{i_p=1}^{n_p} |a_{1, i_1}| \cdots |a_{p, i_p}| N'(f(e_{1, i_1}, \ldots, e_{p, i_p})) \\
        & \leqslant \underbrace{\left( \sum_{i_1=1}^{n_1} \cdots \sum_{i_p=1}^{n_p} N'(f(e_{1, i_1}, \ldots, e_{p, i_p})) \right)}_{ = K \geqslant 0} \prod_{j=1}^p \left( N_{\infty, j}(x_j) \right)
    \end{align*}

    Par le critère de continuité des applications $p$-linéaires, $f$ est continue.
    
\end{proof}

\begin{exemple}
    Soit $E$ un espace préhilbertien réel, muni de $<\cdot, \cdot>$ et de la norme $||\cdot||$ associée.

    Alors $<\cdot, \cdot>$ est continue de $E \times E$ dans $\R$. 

    En munissant $E$ de la norme $||\cdot||_2$. 

    En effet, $\forall (x, y) \in E^2, \, |<x, y>| \leqslant ||x|| \cdot ||y||$ par inégalité de Cauchy-Schwarz,
    et la continuité par le critère. 

    \lvspace Soit $n, p, q \in \N^*$, $\fonction{}{\mathscr M_{n, p}(\K) \times \mathscr M_{p, q}(\K)}{\mathscr M_{n, q}(\K)}{(A, B)}{AB}$ est
    continue car bilinéaire depuis des espaces de dimesnion finie. 


    \lvspace Soient $E, F, G$ espaces vectoriels de dimension finie.
    
    Alors $\fonction{\varphi}{\mathscr L(F, G) \times \mathscr L(E, F)}{\mathscr L(E, G)}{(v, u)}{v \circ u}$ est bilinéaire
    et $||| v \circ u||| \leqslant |||v||| \cdot |||u|||$ donc continue vis à vis des normes d'opérateur par 
    le critère. 
\end{exemple}


\section{Compacts d'un espace vectoriel normé }

\begin{definition}{Compact}{}
    Soit $A \subset E$. 

    On dit que $A$ est compact (ou que $A$ est un compact) ssi toute suite de $A$ possède une 
    valeur d'adhérence dans $A$. 

    C'est à dire pour toute suite $(u_n)_{n \in \N}$ de $A$, on peut en extraire une suite 
    $(u_{\varphi(n)})_{n \in \N}$ telle que $u_{\varphi(n)} \to \ell \in A$.
\end{definition}


\begin{definition}{Borel-Lebesgue}{}
    Soit $A \subset E$, alors
    $A$ est un compact ssi de toute recouvrement de $A$ par des ouverts on peut extraire un sous 
    recouvrement fini, c'est à dire pour toute famille $(U_i)_{i \in I}$ d'ouverts de $E$ 
    telle que $\displaystyle A \subset \bigcup_{i \in I} U_i$, il existe $i_1, \ldots, i_n \in I$ tels que
    $$\displaystyle A \subset \bigcup_{k=1}^n U_{i_k}$$
\end{definition}

\subsection{Propriétés des compacts }

\begin{proposition}{Propriétés des compacts }{}
    \begin{enumerate}
        \item Un compact est fermé et borné.
        \item Un fermé relatif d'un compact est un compact.
        \item Une suite d'un compact converge ssi elle possède une unique valeur d'adhérence. 
        \item Un produit fini de compacts est un compact.
    \end{enumerate}
\end{proposition}

\idee{
    1. caractérisation séquentielle des fermés et absurde 

    2. caractérisation séquentielle 

    3. absurde 

    4. extractions successives
}

\begin{proof}
    1. Soit $A \subset E$ un compact. Montrons que $A$ est fermé.

    Soit $(u_n)_{n \in \N}$ une suite de $A$ qui converge vers $\ell \in E$.

    $A$ étant compact, il existe une sous-suite $(u_{\varphi(n)})_{n \in \N}$ qui converge vers
    $\ell' \in A$. Or $u_{\varphi(n)} \to \ell$ donc $\ell = \ell' \in A$, donc $A$ est fermé.

    Supposons que $A$ n'est pas borné, pour tout $n \in \N$, il existe 
    $x_n \in A$ tel que $N(x_n) > n$.

    $A$ étant compact, on peut extraire de $(x_n)_{n \in \N}$ une sous-suite $(x_{\varphi(n)})_{n \in \N}$ qui converge vers $\ell \in A$.

    Donc $(x_{\varphi(n)})_{n \in \N}$ est bornée, or $\forall n \in \N, \, N(x_{\varphi(n)}) > \varphi(n) \geqslant n$, absurde 

    Donc $A$ est borné.

    \lvspace 2. Soit $V$ un fermé relatif de $A$. Aors $V = \overline V \cap A$. Montrons 
    que $V$ est compact. 

    Soit $(u_n)_{n \in \N}$ une suite de $V$. $A$ étant compact, on peut en extraire 
    une sous-suite $(u_{\varphi(n)})_{n \in \N}$ qui converge vers $\ell \in A$.

    $u_{\varphi(n)} $ est une suite de $V$, donc $\ell \in \overline V$. Donc 
    $\ell \in V$. Donc $V$ est compact.

    \lvspace 3. Toute suite convergente possède une unique valeur d'adhérence d'où $\Rightarrow$. 

    $\Leftarrow$ Soit $(u_n)_{n \in \N}$ une suite d'un compact $A$. Supposons qu'elle possède 
    une unique valeur d'adhérence $a$.

    Supposons par l'absurde que $(u_n)_{n \in \N}$ ne converge pas vers $a$.

    Alors $ \exists \varepsilon > 0, \, \forall N \in \N, \, \exists n \geqslant N, \, d(u_n, a) > \varepsilon$.

    En particulier, $M_\varepsilon = \{ n \in \N, \, d(u_n, a) > \varepsilon \}$ est infini. 

    Numérotons $\varphi(0) < \varphi(1) < \ldots < \varphi(n) < \ldots$ les éléments
    de $M_\varepsilon$. Alors $(u_{\varphi(n)})_{n \in \N}$ est une suite de $A$, donc 
    possède une valeur d'adhérence $b \in A$ qui est aussi une valeur d'adhérence de 
    $(u_n)_{n \in \N}$. Par hypothèse, $b = a$.

    Or $\forall n \in \N, \, N(u_{\varphi(n)} - a) > \varepsilon$, 
    donc à la limite $0 \geqslant \varepsilon$, absurde.

    \lvspace 4. Soient $(E_1, N_1), \ldots, (E_p, N_p)$ des evn tels que 
    $\forall i, \, A_i$ compact de $E_i$ et $A = A_1 \times \cdots \times A_p$. 

    Soit $(u_n)_{n \in \N}$ une suite de $A$.

    Notons $u_n = (u_{n, 1}, \ldots, u_{n, p})$ avec $(u_{n, 1})$ une suite de $A_1$, et $A_1$ est un compact. 

    On peut donc en extraire une sous-suite $(u_{\varphi_1(n), 1})_{n \in \N}$ qui converge vers $\ell_1 \in A_1$.

    De même, on peut extraire de $(u_{\varphi_1(n), 2})_{n \in \N}$ une sous-suite bla bla bla 

    De proches en proches on obtient $u_{\varphi_1 \circ \ldots \circ \varphi_p(n)} \to \ell_p \in A_p$.+

    Donc $(u_{\varphi_1 \circ \ldots \circ \varphi_p(n)}) \to (\ell_1, \ldots, \ell_p) \in A$.

    D'où $A$ compact. 
\end{proof}



\subsection{Cas de la dimension finie }

\begin{theoreme}{Compacts en dimension finie }{}
    Si $(E, N)$ est un evn de dimension finie, tout fermé borné de $E$ est compact. Ainsi, en 
    dimension finie les compacts \textbf{sont} les fermés bornés. 
\end{theoreme}

\idee{
    Bolzano-Weierstrass
}

\begin{proof}
    Soit $A$ un fermé borné de $E$, soit $(u_n)_{n \in \N}$ une suite de $A$.

    Alors $(u_n)$ est bornée donc par Bolzano-Weierstrass, on peut en extraire une
    sous-suite $(u_{\varphi(n)})_{n \in \N}$ qui converge vers $\ell \in E$.

    Or $A$ est fermé donc $\ell \in A$. Donc $A$ est compact.
\end{proof}


\begin{theoreme}{Riesz, HP}{}
    Si $(E, N)$ est de dimension infinie, $B'(0, 1)$ n'est pas compacte. 

    Ainsi un evn est de dimension finie ssi sa boule unité est compacte. 
\end{theoreme}

\begin{proof}
    A la fin du TD.
\end{proof}


\subsection{Compacité et continuité }

\begin{theoreme}{Heine, uniforme continuité}{}
    Toute fonction continue sur un compact est uniformément continue. 
\end{theoreme}

\idee{
    Absurde
}

\begin{proof}
    Soit $A \subset E$ compact et $f \in \mathscr C(A, E')$. 

    Montrons que $f$ est uniformément continue sur $A$. 

    Par l'absurde, supposons que $\exists \varepsilon > 0, \, \forall \eta > 0, \, \exists x, y \in A, \, N(x - y) \leqslant \eta$ et $N'(f(x) - f(y)) > \varepsilon$.

    Soit $n \in \N^*$, $\eta = \frac 1 n$ fournit $x_n, y_n \in A$ tels que $N(x_n - y_n) \leqslant \frac 1 n$
    et $N'(f(x_n) - f(y_n)) > \varepsilon$.

    $A^2$ est compact, donc on peut extraire $(x_{\varphi(n)}, y_{\varphi(n)})_{n \in \N}$ qui converge vers $(\ell_1, \ell_2) \in A^2$.
    $\forall n \in \N, \, N(x_{\varphi(n) - y_{\varphi(n)}}) \leqslant \frac 1 {\varphi(n)}$ et $N'(f(x_{\varphi(n)}) - f(y_{\varphi(n)})) > \varepsilon$.

    En passant à la limite, on obtient $N(\ell_1 - \ell_2) \leqslant 0$ et comme $f$ continue, 
    $N'(f(\ell_1) - f(\ell_2)) \geqslant \varepsilon$.

    Donc $\ell_1 = \ell_2$ et $0 \geqslant \varepsilon$, absurde.
\end{proof}

\begin{theoreme}{Image d'un compact par une application continue}{}
    L'image d'un compact par une application continue est un compact. 
\end{theoreme}

\idee{
    Suites avec fonction
}

\begin{proof}
    Soit $A \subset E$ compact et $f : A \to E'$ continue. Considérons $B = f(A)$. 

    Soit $(b_n)_{n \in \N}$ suite de $B$. Alors $\forall n \in \N, \, \exists a_n \in A, \, b_n = f(a_n)$
    $(a_n)_{n \in \N}$ est une suite de $A$ qui est compact, on peut donc en extraire 
    $a_{\varphi(n)} \limit n \infty \ell \in A$. 

    Donc $b_{\varphi(n)} = f(a_{\varphi(n)}) \limit n \infty f(\ell) $ car $f$ continue en 
    $\ell$ car $f$ continue sur $A$. Or $f(\ell) \in B$, donc $B$ est compact.
\end{proof}

\begin{theoreme}{Théorème des bornes atteintes}{}
    Toute fonction numérique continue sur un compact est bornée et atteint ses bornes, c'est à dire 
    si $A \subset E$ compact et $f \in \mathscr C(A, \R)$, alors 
    $\exists x_1, x_2 \in A$,  $$\forall x \in A, \, f(x_1) \leqslant f(x) \leqslant f(x_2)$$

    Ainsi, $f(x_1) = \min \{ f(x), x \in A \}$ et $f(x_2) = \max \{ f(x), x \in A \}$.
\end{theoreme}

\idee{Caractérisation séquentielle du sup et de l'inf}

\begin{proof}
    Par le théorème précédent, $f(A)$ compact de $\R$, donc borné, $f(A) \neq \emptyset$ donc 
    $ \inf(A)$ et $\sup(A)$ existent par caractérisation séquentielle, et donc il existe une suite 
    $\alpha_n$ de $f(A)$ telle que $\alpha_n \limit n \infty \sup(f(A))$.

    Or $f(A)$ fermé donc $\sup(f(A)) \in f(A)$, donc c'est un max, idem pour l'inf. 
\end{proof}

\subsubsection{Equivalence des normes en dimension finie }

Montrons le théorème d'équivalence des normes en dimension finie.

Démo non exigible.

\idee{
    Montrer que toute norme est équivalente à la norme infinie, et montrer 
    que la norme de $E$ restreinte à la sphère de rayon $1$ et de centre $0$ pour la norme 
    infinie est lipschitzienne. 
}

\begin{proof}
    Soit $E$ un $\K$-ev de dimension finie. 

    Soit $B = (e_1, \ldots, e_p)$ une base de $E$, soit $N$ une norme sur $E$ et $N_\infty$ définie par 
    
    $$N_\infty \left( \sum_{i=1}^p x_i e_i \right) = \max(|x_1|, \ldots, |x_p|)$$

    Considérons $S_\infty (0, 1) = \{ x \in E, \, N_\infty(x) = 1 \}$ et 
    $\fonction f {S_\infty(0, 1)}{\R}{x}{N(x)}$.

    Soient $x, y \in S_\infty(0, 1)$. Alors $| f(x) - f(y)| = |N(x) - N(y)| \leqslant N(x - y)$ par la deuxième inégalité triangulaire. 

    Or \begin{align*}
        N(x - y) & = N\left( \sum_{i=1}^p (x_i - y_i) e_i \right) \\
        & \leqslant \sum_{i=1}^p |x_i - y_i| N(e_i) \\
        & \leqslant \left( \sum_{i=1}^p N_\infty(x - y) N(e_i) \right) \\
        & \leqslant K N_\infty(x - y) \text{ avec } K = \sum_{i=1}^p N(e_i) > 0
    \end{align*}

    Donc $f$ est $K$-lipschitzienne de $(E, N_\infty)$ dans $(\R, |\cdot |)$, donc continue.

    Or $S_\infty(0, 1)$ est fermée bornée de $E$ dans compact car $E$ de dimension finie. 

    Donc (th des bornes atteintes) il existe $x_1$, $x_2 \in S_\infty(0, 1)$ tels que

    $\forall x \in S_\infty(0, 1), \, f(x_1) \leqslant f(x) \leqslant f(x_2)$ et 
    $x_1 \neq 0$ donc $f(x_1) = N(x_1) > 0$. 

    Posons $ \alpha = f(x_1), \, \beta = f(x_2)$. On a $\alpha, \beta > 0$. 

    Soit $y \in E$ tel que $y \neq 0$. Posons $x = \frac y {N_\infty(y)}$ $N_\infty(x) = 1$.

    Donc $\alpha \leqslant f(x) \leqslant \beta$ donc $\alpha N_\infty(y) \leqslant N(y) \leqslant \beta N_\infty(y)$.

    Vrai aussi pour $y = 0$. Donc $N$ et $N_\infty$ sont équivalentes.
\end{proof}


\section{Connexité par arcs }

\begin{definition}{Chemin }{}
    Soit $A \subset E$ non vide, soient $a, b \in A$. 

    On appelle chemin ou arc joignant $a$ à $b$ toute application continue $f$ de 
    $\inter 0 1$ dans $A$ telle que $f(0) = a$ et $f(1) = b$.
\end{definition}


\subsection{Connexité par arcs }


\begin{definition}{Connexité par arcs }{}
    Soit $A \subset E$. 
    On dit que $A$ est connexe par arcs ssi $\forall (a, b) \in A^2, \, \exists f \in \mathscr C(\inter 0 1, A), f(0) = a \text{ et } f(1) = b$. 

    C'est à dire que pour tout $(a, b) \in A^2$, il existe un chemin de $A$ reliant $a$ à $b$.
\end{definition}


\begin{definition}{Ensemble étoilé }{}
    Soit $A \subset E$.
    On dit que $A$ est étoilée par rapport à l'un de ses points ssi $\exists \omega \in A, 
    \, \forall a \in A, \inter a \omega \subset A$.
\end{definition}

\begin{proposition}{Liens entre convexité et étoilé }{}
    Soit $A \subset E$ non vide. 

    Si $A$ est convexe, alors $A$ est étoilée par rapport à tous ses points. 

    Si $A$ est étoilée par rapport à l'un de ses points, $A$ est connexe par arcs. 
\end{proposition}

\idee{poser la bonne fonction }

\begin{proof}
    1. par def de convexe car $A$ non vide. 

    \avspace 2. Supposons que $A$ est étoilée par rapport à $\omega \in A$.

    Alors si $x, y \in A$. Posons $f: \space \begin{array}{|ccc}
    \inter 0 1 & \to & A \\
    t & \mapsto & \left\{ \begin{array}{cc}
        2t \omega + (1 - 2t)x & \text{ si } t \in \inter 0 {\frac 1 2} \\
        (2 - 2t) \omega + (2t - 1)y & \text{ si } t \in \inter {\frac 1 2} 1        
    \end{array} \right.
  \end{array}$

    $f$ est bien continue sur $\interco 0 {\frac 1 2}$ et $\interoc {\frac 1 2} 1$ et 
    $\displaystyle \lim_{t \to \frac 1 2^-} f(t) = \lim_{t \to \frac 1 2^+} f(t) = \omega$ donc $f$ est continue sur $\inter 0 1$.

    $\displaystyle \forall t \in \inter 0 {\frac 1 2}, \, f(t) \in \inter x \omega \subset A$ et
    $\displaystyle \forall t \in \inter {\frac 1 2} 1, \, f(t) \in \inter \omega y \subset A$.

    Donc $f$ est bien à valeurs dans $A$, et $f(0) = x$, $f(1) = y$. 
    Donc $f$ est un chemin de $A$ joignant $x$ à $y$. Donc $A$ est connexe par arcs.
\end{proof}

\subsection{Relation d'équivalence }

\begin{theoreme}{Relation d'équivalence}{}
    Soit $A \subset E$ non vide.

    On définit la relation $\mathcal R_A$ sur $A$ par : 
    $$\forall (a, b) \in A^2, \, a \mathcal R b \Leftrightarrow \exists f \in \mathscr C(\inter 0 1, A), f(0) = a \text{ et } f(1) = b$$

    ou plus simplement $\forall x, y \in A, \, x \mathcal R y$ ssi il existe un chemin de $A$ joignant $x$ à $y$.
    
    Alors $\mathcal R$ est une relation d'équivalence.
\end{theoreme}

\begin{proof}
    \uline{Réflexivité :}

    \svspace Soit $x \in A$. Posons $\fonction f {\inter 0 1} A t {f(t) = x}$. 

    $f$ est continue à valeurs dans $A$, et $f(0) = x$, $f(1) = x$. Donc $x \mathcal R x$.

    \uindent{Symétrie }

    \svspace Soient $x, y \in A$ tels que $x \mathcal R y$. Alors il existe $\fonction f {\inter 0 1} A t {f(t)}$
    telle que $f(0) = x$ et $f(1) = y$.

    Posons $\fonction g {\inter 0 1} A t {g(t) = f(1 - t)}$. On a bien $\forall t \in \inter 0 1, \, g(t) \in A$. 

    $g$ est continue par composition, $g(0) = f(1) = y$ et $g(1) = f(0) = x$. 

    \uindent{Transitivité }

    \svspace Soient $x, y, z \in A$ tels que $x \mathcal R y$ et $y \mathcal R z$.

    Par hypothèse, il existe $f : \inter 0 1 \to A$ continue telle que $f(0) = x$ et $f(1) = y$, 
    et $g : \inter 0 1 \to A$ continue telle que $g(0) = y$ et $g(1) = z$.

    Posons $\fonction h {\inter 0 1} A t { \left\{ \begin{array}{cc} 
        f(2t) & \text{ si } t \leqslant \frac 1 2 \\
        g(2t - 1) & \text{ si } t > \frac 1 2
    \end{array} \right. }$ alors h est bien définie, continue sur $\interco 0 {\frac 1 2}$ et $\interoc {\frac 1 2} 1$ par 
    composition, et $\displaystyle \lim_{t \to \frac 1 2^-} h(t) = \lim_{t \to \frac 1 2^+} h(t) = y$ donc $h$ est continue en $\frac 1 2$.

    De plus $h(0) = f(0) = x$ et $h(1) = g(1) = z$, donc $h$ est bien un chemin de 
    $A$ joignant $x$ à $z$. Donc $x \mathcal R z$.
\end{proof}

\begin{definition}{Composantes connexes par arcs }{}
    Avec ces notations, les classes d'équivalence de $\mathcal R_A$ sont appelées 
    les composantes connexes par arcs de $A$.
\end{definition}

\begin{exemple}
    $A = R^*$ a deux composantes connexes par arcs : ${\R_+}^*$ et ${\R_-}^*$.

    En effet, si $x, y \in {\R_+}^*$, le segment $[x, y] \subset {\R_+}^*$ donc $x \mathcal R y$.
    De même pour ${\R_-}^*$.

    Soient $x \in {\R_+}^*$ et $y \in {\R_-}^*$, et montrons que non $x \mathcal{R} y$.

    S'il existait $f \in \mathscr C(\inter 0 1, \R^*)$ telle que $f(0) = x$ et $f(1) = y$,
    et par le TVI il existerait $t_0 \in \inter 0 1$ tel que $f(t_0) = 0$, absurde.
\end{exemple}

\begin{theoreme}{Image d'une partie connexe par arcs par une application continue }{}
    L'image d'une partie connexe par arcs par une application continue est connexe par arcs.
\end{theoreme}

\begin{proof}
    Soit $f : A \to E'$ continue où $A \subset E$. Supposons que 
    $A$ est connexe par arcs et montrons que $f(A)$ est connexe par arcs. 

    Soient $a, b \in B = f(A)$. Il existe $x, y \in A$ tels que $f(x) = a$ et $f(y) = b$ or 
    $A$ est connexe par arcs. 

    Donc il existe $\varphi : \inter 0 1 \to A$ continue telle que $\varphi(0) = x$ et $\varphi(1) = y$.

    Posons $\fonction \psi {\inter 0 1}{E'} t { f \circ \varphi(t)}$, on a bien 

    \begin{itemize}
        \item $\forall t \in \inter 0 1, \, f \circ \varphi (t) \in B $ 
        \item $\psi$ continue par composition 
        \item $\psi(0) = 0, \, \psi(1) = b$
    \end{itemize}

    Donc $\psi $ est un chemin de $B$ reliant $a$ à $b$. Donc $B$ est connexe par arcs.
\end{proof}

\subsection{Application à \texorpdfstring{$\R$}{R}}

\begin{theoreme}{Connexité par arcs dans R}{}
    Les parties connexes par arcs de $\R$ sont les intervalles.
\end{theoreme}

\begin{proof}
    Soit $A \subset \R$. 

    Si $A$ est un intervalle, alors $A$ est convexe donc $A$ est connexe par arcs. 

    Supposons que $A$ est connexe par arcs. 

    Soient $a \leqslant b \in A$. Montrons que $[a, b] \subset A$.

    Or il existe $\fonction f {\inter 0 1} A t {f(t)}$ continue telle que $f(0) = a$ et $f(1) = b$.

    Soit $c \in \inter a b$. Par le TVI, $f$ étant continue, $\exists t_0 \in \inter 0 1, \, f(t_0) = c$, 
    donc $c \in A$ donc $A$ convexe donc $A$ est un intervalle.
\end{proof}

\begin{theoreme}{Théorème des valeurs intermédiaires }{}
    Soit $A \subset E$ connexe par arcs, $f \in \mathscr C(A, \R)$ et $a, b \in A$.

    Soit $x$ compris entre $f(a)$ et $f(b)$. Alors il existe $c \in A$ tel que $f(c) = x$.
    
\end{theoreme}

\idee{nature image d'un connexe par arcs dans R}

\begin{proof}
    $f(a)$ est connexe par arcs car $A$ l'est et $f$ continue, or $f(a) \subset \R$ donc 
    $f(A)$ est un intervalle, et $f(a), f(b) \in f(A)$ donc $x \in f(A)$. 
\end{proof}

\begin{exemple}
    $A = GL_n(\R)$ alors $\det A \to \R$ est continue, or $det(A) = \R^*$

    Si $\alpha \in \R^*$ posons $M = \operatorname{Diag}(\alpha, 1, \ldots, 1)$. 
    Alors $M \in A$ et $\det(M) = \alpha$.

    Or $\R^*$ n'est pas connexe par arcs donc par contraposée du théorème, $A$ non plus. 
\end{exemple}

\subsection{Complément HP}

\begin{definition}{Connexité HP}{}
    Soit $A \subset E$. 

    On dit que $A$ est connexe ssi toute application $f : A \to \{0, 1\}$ continue est constante. 

\end{definition}

\begin{proposition}{Connexité par arcs implique connexité HP}{}
    Toute partie connexe par arcs est connexe. 
\end{proposition}

\idee{absurde et TVI}

\begin{proof}
    Soit $A \subset E$ connexe par arcs et $f : A \to \{0, 1\}$ continue. Supposons $f$ non constante. 

    Alors $\exists a, b \in A$ tel que $f(a) = 0$ et $f(b) = 1$.

    Il existe $\varphi : \inter 0 1 \to A$ continue telle que $\varphi(0) = a$ et $\varphi(1) = b$.

    Donc $f \circ \varphi : \inter 0 1 \to \{0, 1\}$ est continue, $f \circ \varphi(0) = 0$ et $f \circ \varphi (1) = 1$. 

    Or $\frac 1 2 \in \inter 0 1$, par le TVI, il existe $t \in \inter 0 1$, $f \circ \varphi (t) = \frac 1 2$, absurde. 
    
\end{proof}

\begin{exemple}
    $A \{(t, \sin \frac 1 t), t > 0\} \cup \{(0, u), \, u \in \inter {-1} 1\}$

    On montre que $A$ est connexe mais que $A$ n'est pas connexe par arcs.
\end{exemple}
